\pdfmapfile{+lm.map}
\pdfmapfile{+cm.map}
\pdfmapfile{+symbols.map}
\documentclass[11pt,a4paper]{article}
\usepackage[margin=24mm]{geometry}
\usepackage[T1]{fontenc}
\usepackage{lmodern,amsmath,amssymb,amsthm,mathtools,booktabs,longtable,array,microtype}
\usepackage{xcolor,enumitem,fancyvrb,fvextra}
\usepackage[hidelinks,breaklinks]{hyperref}
\usepackage{bookmark}
\newcommand{\Z}{\mathbb Z}\newcommand{\N}{\mathbb N}
\newcommand{\Con}{\operatorname{Con}}
\newcommand{\tightlist}{\setlength{\itemsep}{0pt}\setlength{\parskip}{0pt}}
\newtheorem{theorem}{Theorem}[section]
\newtheorem{lemma}[theorem]{Lemma}
\theoremstyle{remark}\newtheorem{remark}[theorem]{Remark}
\setlength{\emergencystretch}{3em}
\setlist{itemsep=3pt,topsep=5pt}
\DefineVerbatimEnvironment{verbatim}{Verbatim}{fontsize=\small,breaklines=true,breakanywhere=true,breaksymbolleft={},xleftmargin=1em}
\hypersetup{pdftitle={Shared Pell caps and coprime product equations for ZFC consistency},pdfauthor={Research contribution prepared for Vladimir Reshetnikov}}
\title{Shared Pell caps and coprime product equations for ZFC consistency\\[8pt]\large Round 16 research report and reproducible artifacts}
\author{Prepared for Vladimir Reshetnikov and the ProveIt repository}
\date{25 October 2026}
\begin{document}
\maketitle
\begin{abstract}
One certified Pell ordinate can bound and identify a complete family of requested Pell indices. Composed with a unified factorial and CRT product, this yields quartics in $249$ nonnegative variables, or $278$ with much sparser support. Both improve all four measured size parameters of their corresponding incoming parents. A new signed extension, with cofinal witnesses and delayed power-sign recovery, gives quartics in $324$ unrestricted integers, or $338$ with $4493$ monomials and height $128446$. Separate small-variable constructions reduce arithmetic cost: eleven-integer product equations use $109$ or $110$ binary operations at exact degrees $966$ and $882$, with two-root alternatives $810/112$ and $654/117$; ten-strictly-positive one-square equations use $87$ or $88$ operations at degrees $436$ and $380$. These are domain-specific Pareto tradeoffs, not a single simplicity optimum. Every final equation is a fixed integer polynomial, with literal residuals or complete arithmetic graphs and authentic coefficient payloads. Full new unbounded proofs and reproducible finite checks are supplied. The consistency equivalence retains explicitly identified ordinary source calibrations; this release proves neither consistency nor a new end-to-end proof-assistant endpoint.
\end{abstract}
\clearpage
\tableofcontents
\clearpage
\section{Results and their precise interpretation}
Write $\N=\{0,1,2,\ldots\}$ and $\Z_{>0}=\{1,2,\ldots\}$. These domains are not interchangeable. Every $\Z$ endpoint quantifies all its inputs over unrestricted integers. A named arithmetic register is a polynomial abbreviation, not an additional existential variable. Coefficient height is measured after full collection, and residual rows are not arithmetic-operation counts.

\begin{center}\small
\begin{tabular}{@{}lrrrrr@{}}\toprule
Endpoint and domain&Variables&Degree&Rows&Terms&Height\\\midrule
Pooled natural cap, $\N$&249&4&186&27078&85178741\\
Sparse natural cap, $\N$&278&4&215&5524&83852788\\
Pooled signed cap, $\Z$&324&4&208&5808&128446\\
Sparse signed cap, $\Z$&338&4&222&4493&128446\\\bottomrule
\end{tabular}\end{center}
The incoming natural parent had $261$ variables, $194$ rows, $31679$ monomials and height $85566955$; its sparse companion had $292$, $225$, $6140$ and $83933334$. Each new corresponding endpoint strictly improves all four measures. The signed principal saves $118$ variables relative to round 15's $442$-integer endpoint but increases support from $4652$ to $5808$. Its $338$-integer sparse companion improves variables, rows and support simultaneously while keeping height fixed. These are comparisons within their domains and a fixed direct branch, not universal variable-count records.

\begin{center}\small
\begin{tabular}{@{}lrrr@{}}\toprule
Small-variable endpoint&Domain and inputs&Degree&Operations\\\midrule
Parity-normalized product&$\Z^{11}$&966&109\\
Short-index product&$\Z^{11}$&882&110\\
Nonzero-center two-root&$\Z^{11}$&810&112\\
Weak-scale two-root&$\Z^{11}$&654&117\\
One-square scaled carrier&$\Z_{>0}^{10}$&436&87\\
Lower-degree scaled carrier&$\Z_{>0}^{10}$&380&88\\
\bottomrule
\end{tabular}\end{center}
The integer product pair improves the incoming matched $966/111$ and $882/112$ degree/cost points by two operations each. The positive $436/87$ and $380/88$ points reduce cost below the public $288/89$ point, but increase degree. The two-root integer pair gives further degree/cost tradeoffs, while the retained round-15 $416/125$, $456/124$, $464/119$, and $488/118$ points still offer lower degree. The direct-chart companions at degree/cost $384/92$ and $334/93$ document the same mechanism on a different carrier; they are not claimed to beat the scaled points or the public degree frontier. Operations are live binary additions, subtractions and multiplications, including squares and scalar multiplication; fixed leaves and register reuse are free. Complete expansion of the small-variable polynomials and their coefficient heights is not claimed. Their fixed coefficients are supplied in full, not charged as hidden variables.

\subsection{The logical conclusion and its boundary}
For each equation $P$ on its stated domain $D$, the ordinary proof, formalizable in ZFC with the retained calibration, establishes
\[
 (\exists \boldsymbol{x}\in D^n\;P(\boldsymbol{x})=0)
 \quad\longleftrightarrow\quad
 \text{a finite syntactic ZFC refutation exists}.
\]
Thus the claimed object-theoretic conclusion is
\[
 \mathrm{ZFC}\vdash
 \bigl[\forall\boldsymbol{x}\in D^n\;P(\boldsymbol{x})\ne0\bigr]
 \longleftrightarrow\Con(\mathrm{ZFC}).
\]
This neither asserts that ZFC is consistent nor proves any equation has no zero. The finite checkers verify exact polynomial artifacts and algebraic transformations, not consistency or the unbounded calibration. No new complete Lean, Rocq, or encoded object-ZFC proof object is supplied.

The direct logical interfaces remain the named statements
\begin{quote}\small
\nolinkurl{KEndpoint.c222_iff_zf_refutation}\\
\nolinkurl{ConstructibleChoice.zf_refutation_iff_zfc_refutation}.
\end{quote}
The scalar route retains the source34/native/compiler/refutation calibration and its signed rank, adjacent-Pell and supported-modulus interpolation lemmas. The positive route uses the entire coherent compact44 BFS tuple and its ordinary native/controller/refutation transfer. These coefficient tuples must not be mixed. The archive retains the complete arithmetic proof closure, including \nolinkurl{calculus.tex}, \nolinkurl{primitives.tex}, \nolinkurl{crt.tex}, and \nolinkurl{inherited_onestack_section.tex}; the presence of inherited Lean source does not convert the new arithmetic into a machine-checked endpoint.

\subsection{Which mechanisms are new here}
The incoming Shared Pell Bound construction supplies the unsigned pooled index cap and the idea of a radix chosen after CRT representatives. Round 15 supplies the square-output factorial fixed point and the single factorial/CRT binomial. Their literal composition, horizon-twelve extension, and exact reductions produce the new natural endpoints.

The signed construction extends the cap to unrestricted principal Pell coordinates, proves a complete anchor without unnecessary sign guards, and recovers power signs in an acyclic order. Cofinal horizon and consecutive-product beta-modulus ideas are credited to the incoming Cofinal Witnesses construction. The released signed equation pools ordinate errors only and orients abscissas through actual power requests. The separate general lemma that pools both coordinate signs is proved later for reference but is \emph{not} the literal final compiler.

The eleven-integer product mechanism is inherited from Shared Binomial Certificates; the new parity normalization removes redundant guard scalings. The positive construction adds divisibility by both auxiliary coprimality factors to the strong Pell ordinate, so one square equation recovers the three required roots. The reduced operation counts are measured on complete live graphs, with their higher degree disclosed.

\subsection{Source pins}
The comparison snapshot is the public ProveIt commit
\begin{center}\small\nolinkurl{f4d8db99e060339f3306a04d507c35bc1d0d68c5}.\end{center}
It was committed at 23:45:29 UTC on 24 October 2026, 28 commits after the preceding pin \nolinkurl{45b399ec51dfbecb891dd51d8a9e4b6b7c361ac6}. Eight new incoming archives, including the preceding round's deliverable, were compared. Presence in the incoming directory is not mathematical acceptance or integration into active source. All seven newly acquired byte streams match their exact pinned Git blob identifiers and sizes. Acquisition receipts retain SHA-256 digests and source URLs.

Original coefficient and calibration pins are retained separately in source-binding receipts; a new comparison snapshot does not rewrite their provenance. The report's comparisons are limited to the pinned collection. No literature-wide novelty or global-minimum claim is made.

\section{The fixed polynomial definitions}
For each quartic endpoint the ordered variable array and every residual are supplied in \nolinkurl{equation.json}. If its rows are $R_1,\ldots,R_s$, the equation is exactly
\[
 P(x_0,\ldots,x_{n-1})=\sum_{j=1}^{s}R_j(x_0,\ldots,x_{n-1})^2=0.
\]
No residual is optional. A term \nolinkurl{[c,[i1,...,ik]]} denotes $c\prod_{a=1}^k x_{i_a}$; the empty list denotes a constant. All rows have degree at most two. The independent collected artifact \nolinkurl{expanded_polynomial.json} contains every nonzero coefficient of $P$. This convention, the literal lists, and the explicit variable domains define four ordinary integer polynomials without any oracle or variable-exponent primitive.

The exact endpoint directories are
\begin{quote}\small
\nolinkurl{architecture/unified_cap_h12/natural}\\
\nolinkurl{architecture/unified_cap_sparse_h12/natural}\\
\nolinkurl{direct/integer324}\\
\nolinkurl{direct/integer338}.
\end{quote}
Gate, projection and translation certificates preserve the entire derivation. Rebuilding and independent expansion are separate checks.

The small-variable equations are fully defined by the register displays in their complete proof sections and by complete literal directed acyclic graphs over $+,-,\times$, integer literals, variable leaves, and fixed coefficient leaves. There is no quantified coefficient port. The source34 six-port payload and compact44 four-port payload are supplied as gzip-compressed minimal unsigned little-endian bytes with finite reconstruction recipes. The source34 $M_F$ port is already shifted. Recipe files and notices identify the convention precisely.

\subsection{The two front charts must not be conflated}
The natural composition uses
\[
 \begin{gathered}
 t=t_0+12,\quad R=2+t+x+y+F,\quad K=R,\quad N=B^K,\\
 H=1+C+A_1+\cdots+A_8,\qquad q=(NH)^K.
 \end{gathered}
\]
The literal signed source instead uses
\[
 \begin{gathered}
 t=16+h_0^2,\quad V_0=(F^2+1)^t,\quad y=V_0(V_0-1),\\
 R_0=2+t+x+y+F,\quad K=2R_0,\quad N=B^K,\\
 H_0=1+\sum_{i=1}^9a_i,\qquad q=(BH_0)^K.
 \end{gathered}
\]
Both use $B=\binom{F^2+t}{t}$ only after a noncircular raw-row bootstrap. All powers in these semantic displays are realized by finite polynomial Pell graphs; the literal coefficient artifacts contain no variable-exponent operation. Equality of the desired semantic interfaces is not equality of the two source polynomials or their witness maps.

The following sections reproduce the full new proofs. Literal-register displays use a caret for exponentiation and an asterisk for multiplication. This convention intentionally preserves every formula without Markdown-to-TeX interpretation. Square roots appear only in ordinary proof arguments, never as operations in a final equation.

% Complete unabridged proof from architecture/UNIFIED_CAP_PROOF.md
\section{Unified factorial CRT product and one shared Pell cap}



\subsection{Result and scope}



The principal composed construction has 249 nonnegative integer variables and 186 quadratic residuals. The polynomial is the sum of the squares of every residual in natural/\allowbreak{}equation.json. It has degree four. At the horizon offset 12, its collected polynomial has 27,078 nonzero monomials and coefficient height 85,178,741. This is a composition, not a claim that the supplied shared Pell cap is new: the incoming cap theorem supplies the shared-index mechanism; round15 supplies the square-output factorial/\allowbreak{}product fusion. The additional radix here is (NH)\textasciicircum{}\allowbreak{}K, where H is one plus the sum of the actual independent CRT representatives. It needs neither a representative upper-bound equation nor the factorial numerator in its base. The amplification budget is K=\allowbreak{}R.



The retained C222 calculus calibration and Horn-flow argument are dependencies. With those dependencies the natural zero set is nonempty exactly when the calculus has a finite refutation, and ZFC proves the equivalence between unsolvability and Con(ZFC). This work does not prove consistency and does not supply a new machine-checked object-language ZFC derivation.



All variable powers in this note denote ordinary finite Pell graphs in the emitted arithmetic system. None remains as a primitive in the coefficient artifact. All variables in the natural endpoint range over integers \textgreater{}\allowbreak{}=\allowbreak{}0.



\subsection{Definitions and order of interpretation}



Let t=\allowbreak{}t0+\allowbreak{}h, h\textgreater{}\allowbreak{}=\allowbreak{}12, and let F be the proposed factorial. Put



\begin{verbatim}
V=(F^2+1)^t, B=choose(F^2+t,t), R=2+t+x+y+F,
K=R, N=B^K, H=1+C+A1+...+A8, q=(NH)^K.
\end{verbatim}

The notation B=\allowbreak{}choose(...) is initially only a proposed binomial output. The raw rows are



\begin{verbatim}
V+B=BF+u+1, B=u+v+1.
\end{verbatim}

The nine representatives C,A1,...,A8 are positive by their raw affine offsets. Therefore H\textgreater{}\allowbreak{}=\allowbreak{}10. The raw shared-cap power theorem first applies to every independently positive base. In particular it gives V=\allowbreak{}(F\textasciicircum{}\allowbreak{}2+\allowbreak{}1)\textasciicircum{}\allowbreak{}t\textgreater{}\allowbreak{}=\allowbreak{}1. The second raw row gives B\textgreater{}\allowbreak{}=\allowbreak{}1 and subtracting the rows gives V+\allowbreak{}v=\allowbreak{}BF. Thus F\textgreater{}\allowbreak{}=\allowbreak{}1, V\textgreater{}\allowbreak{}F\textasciicircum{}\allowbreak{}2 and BF\textgreater{}\allowbreak{}=\allowbreak{}V, whence B\textgreater{}\allowbreak{}F and V\textless{}\allowbreak{}B\textasciicircum{}\allowbreak{}2. The positive-base theorem now gives N=\allowbreak{}B\textasciicircum{}\allowbreak{}K and q=\allowbreak{}(NH)\textasciicircum{}\allowbreak{}K. Since K\textgreater{}\allowbreak{}=\allowbreak{}2, N\textgreater{}\allowbreak{}=\allowbreak{}B\textasciicircum{}\allowbreak{}2\textgreater{}\allowbreak{}V and q\textgreater{}\allowbreak{}=\allowbreak{}N\textgreater{}\allowbreak{}V. Moreover N divides q. This argument uses no binomial semantics, decoded representative, or factorial identity, and hence is noncircular.



For positive r and n, every coefficient choose(r+\allowbreak{}n,j), 0\textless{}\allowbreak{}=\allowbreak{}j\textless{}\allowbreak{}=\allowbreak{}n, is at most (r+\allowbreak{}1)\textasciicircum{}\allowbreak{}n. The proof is induction using Pascal's identity for j\textless{}\allowbreak{}n and the product formula for j=\allowbreak{}n. Applying it with r=\allowbreak{}F\textasciicircum{}\allowbreak{}2,n=\allowbreak{}t shows that all coefficients relevant to B are \textless{}\allowbreak{}=\allowbreak{}V\textless{}\allowbreak{}q. The exact digit graph therefore really gives B=\allowbreak{}choose(F\textasciicircum{}\allowbreak{}2+\allowbreak{}t,t).



\subsection{The square factorial fixed point at horizon twelve}



Write T=\allowbreak{}t! and f=\allowbreak{}V/\allowbreak{}B. The two ceiling rows are equivalent to F-1\textless{}\allowbreak{}f\textless{}\allowbreak{}=\allowbreak{}F, so F=\allowbreak{}ceil(f). Also



\begin{verbatim}
f=T product_(i=1)^t (F^2+1)/(F^2+i).
\end{verbatim}

Every factor is \textless{}\allowbreak{}=\allowbreak{}1, hence F\textless{}\allowbreak{}=\allowbreak{}T. Each factor is nondecreasing in F\textasciicircum{}\allowbreak{}2; since F\textgreater{}\allowbreak{}=\allowbreak{}1, f\textgreater{}\allowbreak{}=\allowbreak{}2\textasciicircum{}\allowbreak{}t/\allowbreak{}(t+\allowbreak{}1)\textgreater{}\allowbreak{}t\textasciicircum{}\allowbreak{}2 for t\textgreater{}\allowbreak{}=\allowbreak{}12. The last inequality starts from 4096\textgreater{}\allowbreak{}12\textasciicircum{}\allowbreak{}2\textasteriskcentered{}\allowbreak{}13=\allowbreak{}1872; its induction step uses t\textasciicircum{}\allowbreak{}2\textgreater{}\allowbreak{}=\allowbreak{}3t+\allowbreak{}2. Thus F\textgreater{}\allowbreak{}t\textasciicircum{}\allowbreak{}2. For ai=\allowbreak{}(i-1)/\allowbreak{}(F\textasciicircum{}\allowbreak{}2+\allowbreak{}i),



\begin{verbatim}
sum ai<t(t-1)/(2F^2)<1/(2F).
\end{verbatim}

The finite product bound product(1-ai)\textgreater{}\allowbreak{}=\allowbreak{}1-sum ai gives T(1-1/\allowbreak{}(2F))\textless{}\allowbreak{}f\textless{}\allowbreak{}=\allowbreak{}F. Consequently



\begin{verbatim}
T<2F^2/(2F-1)=F+F/(2F-1)<=F+1.
\end{verbatim}

Integrality gives T\textless{}\allowbreak{}=\allowbreak{}F, so F=\allowbreak{}t!. Conversely for F=\allowbreak{}T the same estimate gives T-1/\allowbreak{}2\textless{}\allowbreak{}f\textless{}\allowbreak{}=\allowbreak{}T. The explicit natural ceiling witnesses therefore exist. This extends round15's stated horizon-16 proof to horizon 12, without changing its identity or introducing a new fixed point.



\subsection{CRT product and the self-sized radix}



For 1\textless{}\allowbreak{}=\allowbreak{}i\textless{}\allowbreak{}=\allowbreak{}t set di=\allowbreak{}1+\allowbreak{}F\textasciicircum{}\allowbreak{}2/\allowbreak{}i. Since F=\allowbreak{}t!, every di is an integer congruent to one modulo F. Every prime divisor of di is greater than t. If a prime divided di and dj, it would divide i-j, so the di are pairwise coprime. Their product is exactly B=\allowbreak{}choose(F\textasciicircum{}\allowbreak{}2+\allowbreak{}t,t).



Every independent representative n is strictly less than H. The lower index of each of its range graphs is t, t+\allowbreak{}1, x+\allowbreak{}1 or y+\allowbreak{}1, and is at most R-1\textless{}\allowbreak{}K. For every 0\textless{}\allowbreak{}=\allowbreak{}j\textless{}\allowbreak{}=\allowbreak{}k,



\begin{verbatim}
choose(n,j)<=max(1,n^k)<H^K<=q.
\end{verbatim}

Coefficients with j\textgreater{}\allowbreak{}n are zero. Thus these digit graphs recover exactly the nine required divisibilities by N. This proof imposes no bound on the actual representative and deletes the round15 aggregate-cap row and slack.



Let J be the derived nonnegative label sum. The retained literal majorant certificate gives J(A)\textless{}\allowbreak{}=\allowbreak{} (n+\allowbreak{}1)\textasciicircum{}\allowbreak{}ell with ell=\allowbreak{}4,194,304 and n the largest of the four encoded schema parameters. Since n+\allowbreak{}1\textless{}\allowbreak{}=\allowbreak{}H, J(A)\textless{}\allowbreak{}=\allowbreak{}H\textasciicircum{}\allowbreak{}ell. At t=\allowbreak{}12, 3ell(t+\allowbreak{}1)=\allowbreak{}163,577,856\textless{}\allowbreak{}958,003,200=\allowbreak{}2\textasteriskcentered{}\allowbreak{}12!. The inequality persists under increasing t because factorial growth dominates t+\allowbreak{}1. Hence



\begin{verbatim}
ell(t+1)<=2F/3<F<K.
\end{verbatim}

All coefficients through index t+\allowbreak{}1 in the derived-label digit are consequently \textless{}\allowbreak{}q. That graph therefore gives N dividing choose(J(A),t+\allowbreak{}1), without putting J(A) in H or requiring an upper bound on its encoded value.



The retained local nonnegative source majorant has E=\allowbreak{}8,388,608. For t\textgreater{}\allowbreak{}=\allowbreak{}12, F\textgreater{}\allowbreak{}=\allowbreak{}12!\textgreater{}\allowbreak{}2\textasciicircum{}\allowbreak{}28, 28E=\allowbreak{}234,881,024\textless{}\allowbreak{}2\textasciicircum{}\allowbreak{}28 and 2E\textless{}\allowbreak{}=\allowbreak{}2\textasciicircum{}\allowbreak{}28. At r=\allowbreak{}2\textasciicircum{}\allowbreak{}28 this proves r\textasciicircum{}\allowbreak{}E\textless{}\allowbreak{}2\textasciicircum{}\allowbreak{}r; for successive r, (1+\allowbreak{}1/\allowbreak{}r)\textasciicircum{}\allowbreak{}E\textless{}\allowbreak{}2 follows by bounding the binomial sum by a geometric series with ratio E/\allowbreak{}r\textless{}\allowbreak{}=\allowbreak{}1/\allowbreak{}2. Thus every local source value and local label sum is \textless{}\allowbreak{}R\textasciicircum{}\allowbreak{}E\textless{}\allowbreak{}2\textasciicircum{}\allowbreak{}R=\allowbreak{}2\textasciicircum{}\allowbreak{}K. This is the required no-wrap bound with K=\allowbreak{}R, not 2R.



\subsection{CRT soundness and completeness}



The inherited prime-power binomial lemma says that p\textasciicircum{}\allowbreak{}e dividing choose(a,k) forces a modulo p\textasciicircum{}\allowbreak{}e to lie in [0,k-1]. For row i-1 choose p dividing di and m=\allowbreak{}vp(di). The raw B divisibility of C+\allowbreak{}F\textasciicircum{}\allowbreak{}2+\allowbreak{}1 identifies C modulo p with i-1. The index range then identifies its small residue modulo p\textasciicircum{}\allowbreak{}(Km) with exactly that index. The other ranges decode every local coordinate in its prescribed interval. The local divisibility by N and the preceding local bound force the nonnegative local source value to be zero, not merely congruent to zero. The derived-label range and its local bound likewise force J(local)\textless{}\allowbreak{}=\allowbreak{}t. The unchanged beta consistency and pairwise-coprime weight flow then recover a finite C222 refutation.



Conversely start with a finite refutation, expand to a proof tree, pad the horizon beyond all required parameters and label sums, and choose t\textgreater{}\allowbreak{}=\allowbreak{}h. The inherited beta construction gives x,y and the local source rows. Set F=\allowbreak{}t!, then B,K,N. Choose positive independent representatives by CRT modulo N\textasciicircum{}\allowbreak{}2. There is no restriction on their size. For each range, the precision-2K congruence supplies at least 2Km prime-adic digits in the falling-factorial numerator, while division by k! loses at most k-1\textless{}\allowbreak{}K\textless{}\allowbreak{}=\allowbreak{}Km. Thus N divides every required binomial. The same proof applies to the derived polynomial J because polynomial congruences preserve the local label. Choose H and q afterwards. The carry bounds apply regardless of these representatives' sizes; take Euclidean digit witnesses and the simultaneous positive-power witnesses. This constructs every raw witness without a fixed-point dependence on the radix.



\subsection{Shared Pell cap interface and logical independence}



The composed source retains exactly the incoming pooled cap theorem. Its W is 2 plus the sum of distinct actual power bases, exponent atoms and outputs. There are 15 distinct positive independent index atoms and one additive index. Each independent k is \textless{}\allowbreak{}W, and their number is \textless{}\allowbreak{}=\allowbreak{}W-2. The scale equation forces A-1\textgreater{}\allowbreak{}W\textasciicircum{}\allowbreak{}3 before any power has been interpreted. One complete indexed-Pell certificate gives Vcap=\allowbreak{}Y\_\allowbreak{}W(A). The aggregate row



\begin{verbatim}
Vcap=W+1+s+sum(Y_i)
\end{verbatim}

bounds all independent principal ordinates; the pooled congruence is



\begin{verbatim}
sum(Y_i-k_i)^2=(A-1)Q.
\end{verbatim}

Positive Pell classification gives Y\_\allowbreak{}i=\allowbreak{}Y\_\allowbreak{}ji(A) with 0\textless{}\allowbreak{}=\allowbreak{}ji\textless{}\allowbreak{}W. Reduction modulo A-1 identifies the pooled error with sum(ji-ki)\textasciicircum{}\allowbreak{}2, lying in [0,W\textasciicircum{}\allowbreak{}3). Thus every ji=\allowbreak{}ki. All sign/\allowbreak{}domain assumptions here are the natural ones. The incoming cap theorem gives constructive completeness, including a nonnegative Q and the aggregate slack. The single additive index is supplied by exact Pell addition, not by an unchecked seed. Power congruences then recover every requested exponentiation. The cap witnesses do not enter the semantic frame.



The raw source is therefore equivalent to the retained source theorem. The exact natural-domain contraction restores every eliminated coordinate and preserves existence in both directions. Finite expansion and regression tests verify emitted syntax and examples, not the unbounded source theorem itself.



\subsection{Alternate radices}



The builder also emits q=\allowbreak{}(BH)\textasciicircum{}\allowbreak{}K and q=\allowbreak{}N H\textasciicircum{}\allowbreak{}K. Both are multiples of N=\allowbreak{}B\textasciicircum{}\allowbreak{}K, both are \textgreater{}\allowbreak{}=\allowbreak{}N\textgreater{}\allowbreak{}V and \textgreater{}\allowbreak{}H\textasciicircum{}\allowbreak{}K, so all arguments above remain valid. At offset 12 they also have 249 natural variables and 186 rows; the principal (NH)\textasciicircum{}\allowbreak{}K endpoint is less dense. These alternates are recorded as measured comparisons, not additional logical assumptions.



\subsection{Sparse companion and the density tradeoff}



The sparse companion retains one inequality Y\_\allowbreak{}i\textless{}\allowbreak{}Vcap and one congruence Y\_\allowbreak{}i-ki=\allowbreak{}(A-1)qi for each independent seed, rather than the aggregate and pooled square rows. It also retains Vcap=\allowbreak{}W+\allowbreak{}margin. Positive Pell classification and the cap give 0\textless{}\allowbreak{}=\allowbreak{}ji\textless{}\allowbreak{}W, while the individual congruence gives ji=\allowbreak{}ki modulo A-1. Since both indices lie in [0,W-1] and W\textless{}\allowbreak{}A-1, they are equal. Canonical pairs supply the nonnegative individual quotients and slacks, so completeness is unchanged. This is the supplied incoming sparse theorem composed with exactly the same fused factorial and pure self-sized radix.



At offset 12 the sparse endpoint has 278 nonnegative variables, 215 quadratic rows, 5,524 quartic monomials and height 83,852,788. The incoming sparse endpoint had 292 variables, 225 rows, 6,140 monomials and height 83,933,334. The new sparse endpoint improves each of those four measures. The 249-variable principal endpoint saves a further 29 variables but has 27,078 monomials. These are two Pareto alternatives; lower variable count does not mean lower density.

% Complete unabridged proof from direct/PROOF.md
\section{Signed pooled Pell bounds and cofinal direct certificates}



\subsection{Fixed equation and dependencies}



The endpoint in \texttt{integer324/\allowbreak{}equation.json} is an ordinary sum of squares of 208 explicitly listed quadratic integer residuals in 324 unrestricted integer unknowns. The companion \texttt{integer338/\allowbreak{}equation.json} replaces one pooled congruence by fifteen individual congruences. Its 338 unknowns and 222 quadratic residuals trade variables for a smaller collected polynomial. The 324-variable polynomial has 5,808 collected monomials; the 338-variable polynomial has 4,493. Both have exact degree four and coefficient height 128,446. The exact collected coefficients are in each endpoint's \texttt{expanded\_\allowbreak{}polynomial.json}. No variable exponent is an operation of either polynomial.



The source is the pinned round15 direct square-shell, derived-label, unified-factorial/\allowbreak{}product construction. Its 18 power requests have 15 independent indices and one derived index. The natural arithmetic and numerical C222 calibration remain the dependencies described in round15 \texttt{architecture/\allowbreak{}UNIFIED\_\allowbreak{}PRODUCT\_\allowbreak{}PROOF.md} and \texttt{direct/\allowbreak{}PROOF.md}. In particular the named logical interfaces remain \texttt{KEndpoint.c222\_\allowbreak{}iff\_\allowbreak{}zf\_\allowbreak{}refutation} and \texttt{ConstructibleChoice.zf\_\allowbreak{}refutation\_\allowbreak{}iff\_\allowbreak{}zfc\_\allowbreak{}refutation}.



This note proves a different signed representation of the complete arithmetic interface, plus an existence-preserving cofinal restriction and a self-sized radix. It proves, in ordinary mathematics formalizable in ZFC, that an integer zero exists exactly when a finite ZFC refutation exists. Negation gives the ZFC-provable equivalence of unsolvability with Con(ZFC). It neither proves consistency nor supplies a new end-to-end proof-assistant derivation. No global minimality or literature-wide novelty is claimed.



The shared-cap idea is inherited from the incoming Shared Pell Bound 261 package. The signed sum-of-squares cap, signed-index recovery, delayed front signs, odd private parameters, and exact composition here are the new signed extension. The cofinal horizon and consecutive-product beta modulus are inherited from the incoming Cofinal Witnesses package, transferred here to the already existing square-output factorial numerator. The self-sized radix is also developed in this round's natural-number construction; the literal signed source here uses q=\allowbreak{}(BH)\textasciicircum{}\allowbreak{}K, not that construction's q=\allowbreak{}(NH)\textasciicircum{}\allowbreak{}K. Both variants have the necessary bounds.



\subsection{The literal semantic changes}



All notation in this section names intended values or retained coordinates. Powers are interpreted only in the order proved below.



The retained front has



\begin{verbatim}
t = 16+t0,       rho=F^2,       V0=(rho+1)^t,
B=binom(rho+t,t),       R0=2+t+x+y+F,       K=2R0,
N=B^K.
\end{verbatim}

The cofinal source adds t0=\allowbreak{}h0\textasciicircum{}\allowbreak{}2 and y=\allowbreak{}V0(V0-1). Let a\_\allowbreak{}1,...,a\_\allowbreak{}9 denote the independent positive CRT representatives, including the encoded index, and set



\begin{verbatim}
H0=1+sum_i a_i,       q=(B H0)^K.
\end{verbatim}

The aggregate constraint sum\_\allowbreak{}i a\_\allowbreak{}i+\allowbreak{}slack=\allowbreak{}9N\textasciicircum{}\allowbreak{}2 is removed. Its only purpose was the old radix bound. The semantic base of the third power request is replaced by B H0. The explicit transformation lists the three changed private power-base rows and the new quadratic product gate.



The old per-index comparison norms, comparison congruences, private ordinate inequalities, and common product over all principal ordinates are removed. Each original principal norm is retained, at its original undoubled requested index k\_\allowbreak{}i. All original power congruences and front remainders are retained without squaring outputs or residue-one restrictions.



The private square frame is



\begin{verbatim}
W=2+sum(a^2 : a is a distinct actual base, exponent, or output atom).
\end{verbatim}

All atoms include their complete affine restorations. Consequently |a|\textless{}\allowbreak{}W, a\textasciicircum{}\allowbreak{}2\textless{}\allowbreak{}W and W\textgreater{}\allowbreak{}=\allowbreak{}258 because t\textgreater{}\allowbreak{}=\allowbreak{}16. For each original requested index, W\textgreater{}\allowbreak{}=\allowbreak{}2|k\_\allowbreak{}i|+\allowbreak{}1. There are m=\allowbreak{}15 principal pairs, so m\textless{}\allowbreak{}W. No positivity of an unvalidated output or of K is used in these facts.



The old frame norm is retained, with its ordinate parameter replaced by an odd integer:



\begin{verbatim}
A^2-W^3(W+2)(2z+1)^2=1,       A>=0.
\end{verbatim}

There is no sign guard on z. The source still has an ordinary quadratic multiplication circuit for this displayed higher-degree identity.



\subsection{Frame and Pell facts}



For integer a\textgreater{}\allowbreak{}=\allowbreak{}2 let X\_\allowbreak{}j(a),Y\_\allowbreak{}j(a) be the Pell sequences with X\_\allowbreak{}0=\allowbreak{}1,Y\_\allowbreak{}0=\allowbreak{}0 and



\begin{verbatim}
X_(j+1)=a X_j+(a^2-1)Y_j,    Y_(j+1)=X_j+a Y_j.
\end{verbatim}

Every integral norm-one solution at discriminant a\textasciicircum{}\allowbreak{}2-1 has independent signs on X\_\allowbreak{}j(a),Y\_\allowbreak{}j(a), for a unique j\textgreater{}\allowbreak{}=\allowbreak{}0. This follows by elementary descent by the smallest positive unit a+\allowbreak{}sqrt(a\textasciicircum{}\allowbreak{}2-1). In particular Y\_\allowbreak{}j\textgreater{}\allowbreak{}=\allowbreak{}j, and the positive ordinate sequence is strictly increasing.



The frame norm and 2z+\allowbreak{}1 nonzero give A\textgreater{}\allowbreak{}=\allowbreak{}2. Apply classification at center W+\allowbreak{}1 to (A,|W(2z+\allowbreak{}1)|). Since Y\_\allowbreak{}j(W+\allowbreak{}1)=\allowbreak{}j modulo W, the positive frame index is divisible by W and is at least W. Hence



\begin{verbatim}
A>=X_W(W+1)>W^W,       A-1>4W^3.
\end{verbatim}

The last inequality holds for W\textgreater{}\allowbreak{}=\allowbreak{}258. All following applications use A positive.



We use two elementary inherited index facts. If Y\_\allowbreak{}i(A)\textasciicircum{}\allowbreak{}2 divides Y\_\allowbreak{}l(A), with i,l\textgreater{}\allowbreak{}=\allowbreak{}1, then i divides l and Y\_\allowbreak{}i(A) divides l. To see this, divide l by i, use the addition law modulo Y\_\allowbreak{}i to force zero remainder, and expand modulo Y\_\allowbreak{}i\textasciicircum{}\allowbreak{}2. Also, for 1\textless{}\allowbreak{}=\allowbreak{}i\textless{}\allowbreak{}=\allowbreak{}l, the congruence X\_\allowbreak{}j(A)=\allowbreak{}+\allowbreak{}/\allowbreak{}-X\_\allowbreak{}i(A) modulo X\_\allowbreak{}l(A) forces j=\allowbreak{}+\allowbreak{}/\allowbreak{}-i modulo 2l. Reduce the Pell unit exponent modulo 2l, use its value -1 at exponent 2l modulo X\_\allowbreak{}l, and reflect about l; strict growth and 2X\_\allowbreak{}(l-1)\textless{}\allowbreak{}X\_\allowbreak{}l distinguish all remaining absolute residues.



These are unbounded arithmetic arguments, not consequences of finite samples.



\subsection{One fully certified anchor}



Introduce the single anchor (U,C), its common pair (H,J), comparison pair (S,Q), and comparison center D. The literal quadratic circuit represents



\begin{verbatim}
U^2-(A^2-1)C^2=1,
H^2-(A^2-1)J^2=1,
S^2-(D^2-1)Q^2=1,
J=C^2(e+1),           e>=0,
D=1+4J(2d+1),
D=A+Hf,              S=U+Hg,              Q=W+4Ch.
\end{verbatim}

All letters other than the stated e are unrestricted integers. The cap row is



\begin{verbatim}
C=W+1+sum_i y_i^2+u^2+v^2+w(w+1).                 (1)
\end{verbatim}

Thus C\textgreater{}\allowbreak{}W\textgreater{}\allowbreak{}=\allowbreak{}258. The product gives J\textgreater{}\allowbreak{}0; hence |H|\textgreater{}\allowbreak{}=\allowbreak{}A and |D|\textgreater{}\allowbreak{}=\allowbreak{}3. D is odd and D=\allowbreak{}1 modulo 4C. No sign of D or d is assumed.



Write |U|=\allowbreak{}X\_\allowbreak{}i(A), C=\allowbreak{}Y\_\allowbreak{}i(A), |H|=\allowbreak{}X\_\allowbreak{}l(A), J=\allowbreak{}Y\_\allowbreak{}l(A), with positive i,l. Squared divisibility gives i\textless{}\allowbreak{}=\allowbreak{}l and C divides l. Classify the comparison norm at |D|. Polynomial base transfer D=\allowbreak{}A modulo H and S=\allowbreak{}U modulo H give comparison index j=\allowbreak{}+\allowbreak{}/\allowbreak{}-i modulo 2l; signs due to a negative D are absorbed by the independent coordinate signs. Since D=\allowbreak{}1 modulo 4C, the ordinate polynomial gives j=\allowbreak{}+\allowbreak{}/\allowbreak{}-W modulo 4C. It follows that i=\allowbreak{}+\allowbreak{}/\allowbreak{}-W modulo 2C. But 1\textless{}\allowbreak{}=\allowbreak{}i\textless{}\allowbreak{}=\allowbreak{}C and W\textless{}\allowbreak{}C, so only i=\allowbreak{}W is possible. Therefore



\begin{verbatim}
C=Y_W(A).
\end{verbatim}

The argument does not need a sign guard on U,H,S or Q. Unlike the old individual certificate, this complete index verification is performed only once.



\subsection{Bounded signed principal indices}



Each retained norm gives



\begin{verbatim}
x_i=sigma_i X_(j_i)(A),       y_i=epsilon_i Y_(j_i)(A),
\end{verbatim}

where j\_\allowbreak{}i\textgreater{}\allowbreak{}=\allowbreak{}0 and sigma\_\allowbreak{}i,epsilon\_\allowbreak{}i belong to \{+\allowbreak{}1,-1\}. Equation (1) implies |y\_\allowbreak{}i|\textless{}\allowbreak{}C, hence j\_\allowbreak{}i\textless{}\allowbreak{}W. Reduce the Pell recurrence modulo A-1 to obtain X\_\allowbreak{}j=\allowbreak{}1 and Y\_\allowbreak{}j=\allowbreak{}j modulo A-1.



The primary source imposes the single quadratic row



\begin{verbatim}
sum_i (y_i-k_i)^2=(A-1)q0.                       (2)
\end{verbatim}

The quotient q0 is unrestricted. Consequently A-1 divides



\begin{verbatim}
E=sum_i (epsilon_i j_i-k_i)^2.
\end{verbatim}

Because j\_\allowbreak{}i\textless{}\allowbreak{}W, |k\_\allowbreak{}i|\textless{}\allowbreak{}W and m\textless{}\allowbreak{}W, one has 0\textless{}\allowbreak{}=\allowbreak{}E\textless{}\allowbreak{}4mW\textasciicircum{}\allowbreak{}2\textless{}\allowbreak{}4W\textasciicircum{}\allowbreak{}3\textless{}\allowbreak{}A-1. Thus E=\allowbreak{}0, and epsilon\_\allowbreak{}i j\_\allowbreak{}i=\allowbreak{}k\_\allowbreak{}i for each i. This statement allows zero or negative proposed k\_\allowbreak{}i: the conclusion is y\_\allowbreak{}i=\allowbreak{}Y\_\allowbreak{}(k\_\allowbreak{}i)(A) in the signed-index convention Y\_\allowbreak{}(-k)=\allowbreak{}-Y\_\allowbreak{}k. No premature positivity of F, K or y is required.



The sparse source instead imposes fifteen rows



\begin{verbatim}
y_i-k_i=(A-1)q_i.
\end{verbatim}

The signed error epsilon\_\allowbreak{}i j\_\allowbreak{}i-k\_\allowbreak{}i has absolute value below 2W\textless{}\allowbreak{}A-1, so it vanishes directly. The same ordinate conclusion follows. The primary and sparse sources therefore have identical semantic interfaces.



The conservative \texttt{cap\_\allowbreak{}both} checkpoint additionally pools (x\_\allowbreak{}i-1)\textasciicircum{}\allowbreak{}2; its residual small error includes (sigma\_\allowbreak{}i-1)\textasciicircum{}\allowbreak{}2 and forces sigma\_\allowbreak{}i=\allowbreak{}+\allowbreak{}1 immediately. The smaller-support released variants omit those terms and orient abscissas by their genuine power requests below.



\subsection{Acyclic power and sign recovery}



The only front sign premises needed at the start are x\textgreater{}\allowbreak{}=\allowbreak{}0, the nine encoded representatives minus one \textgreater{}\allowbreak{}=\allowbreak{}0, factorial remainder r\textgreater{}\allowbreak{}=\allowbreak{}0, factorial slack s\textgreater{}\allowbreak{}=\allowbreak{}0, and the exact-binomial output slack o\textgreater{}\allowbreak{}=\allowbreak{}0. The literal raw rows give



\begin{verbatim}
B=r+s+1>=1,          q=B+o+1>=2.
\end{verbatim}

The cofinal horizon gives t\textgreater{}\allowbreak{}=\allowbreak{}16. Thus the pooled conclusion already gives the t ordinate Y\_\allowbreak{}t(A)\textgreater{}\allowbreak{}0, although its abscissa might still have either sign.



For any positive base 2\textless{}\allowbreak{}=\allowbreak{}n\textless{}\allowbreak{}W, exponent 1\textless{}\allowbreak{}=\allowbreak{}k\textless{}\allowbreak{}W and candidate output |m|\textless{}\allowbreak{}W, put T=\allowbreak{}2An-n\textasciicircum{}\allowbreak{}2-1. Then T\textgreater{}\allowbreak{}W\textasciicircum{}\allowbreak{}W+\allowbreak{}W. If the abscissa is positive, the power congruence gives m=\allowbreak{}n\textasciicircum{}\allowbreak{}k modulo T and the size bounds give equality. If it is negative while the ordinate is positive, the Pell norm instead gives m\textasteriskcentered{}\allowbreak{}n\textasciicircum{}\allowbreak{}k=\allowbreak{}-1 modulo T. The absolute value of m\textasteriskcentered{}\allowbreak{}n\textasciicircum{}\allowbreak{}k+\allowbreak{}1 is below T, so equality m\textasteriskcentered{}\allowbreak{}n\textasciicircum{}\allowbreak{}k=\allowbreak{}-1 would follow, impossible for n\textgreater{}\allowbreak{}=\allowbreak{}2. This ordinary signed-output lemma uses no sign of m.



Apply it first to the genuine request q\textasciicircum{}\allowbreak{}t. It orients the t abscissa positively. The same positive pair now certifies V0=\allowbreak{}(F\textasciicircum{}\allowbreak{}2+\allowbreak{}1)\textasciicircum{}\allowbreak{}t, including the possible base one, by direct small-residue uniqueness. The ceiling identity is FB=\allowbreak{}V0+\allowbreak{}s\textgreater{}\allowbreak{}0. Since B\textgreater{}\allowbreak{}0, F\textgreater{}\allowbreak{}=\allowbreak{}1. Also V0\textgreater{}\allowbreak{}=\allowbreak{}F\textasciicircum{}\allowbreak{}2+\allowbreak{}1 and FB\textgreater{}\allowbreak{}=\allowbreak{}V0, so B\textgreater{}\allowbreak{}F\textgreater{}\allowbreak{}=\allowbreak{}1, B\textgreater{}\allowbreak{}=\allowbreak{}2 and q\textgreater{}\allowbreak{}=\allowbreak{}3. None of these conclusions uses a binomial interpretation.



Now y=\allowbreak{}V0(V0-1)\textgreater{}\allowbreak{}=\allowbreak{}F+\allowbreak{}1: already V0\textgreater{}\allowbreak{}=\allowbreak{}F\textasciicircum{}\allowbreak{}2+\allowbreak{}1, and if F=\allowbreak{}1, t\textgreater{}\allowbreak{}=\allowbreak{}16 gives V0=\allowbreak{}2\textasciicircum{}\allowbreak{}t\textgreater{}\allowbreak{}=\allowbreak{}3; if F\textgreater{}\allowbreak{}=\allowbreak{}2, V0\textgreater{}\allowbreak{}=\allowbreak{}F+\allowbreak{}2. Hence the deleted endpoint slack r0=\allowbreak{}y-F-1 is nonnegative. R0=\allowbreak{}2+\allowbreak{}t+\allowbreak{}x+\allowbreak{}y+\allowbreak{}F and K=\allowbreak{}2R0 are positive. Every representative exponent is positive from its guarded +\allowbreak{}1 constructor, and the derived label exponent is positive by the inherited monotone constructor circuit. Thus every requested index is now positive.



The B\textasciicircum{}\allowbreak{}K request has base B\textgreater{}\allowbreak{}=\allowbreak{}2 and orients the K pair, giving N=\allowbreak{}B\textasciicircum{}\allowbreak{}K. Its other request has base B H0\textgreater{}\allowbreak{}=\allowbreak{}2 and gives q=\allowbreak{}(B H0)\textasciicircum{}\allowbreak{}K. Every remaining principal pair supports a request at base q or q+\allowbreak{}1, so the same lemma orients it and certifies its power. The unchanged one-step Pell addition from the oriented t pair certifies the derived t+\allowbreak{}1 request. All eighteen powers have now been recovered.



For every exact or range digit, replace its two old remainder-sign guards by the one guard r\_\allowbreak{}i\textasteriskcentered{}\allowbreak{}s\_\allowbreak{}i\textgreater{}\allowbreak{}=\allowbreak{}0. Its unchanged bound row is L\_\allowbreak{}i=\allowbreak{}r\_\allowbreak{}i+\allowbreak{}s\_\allowbreak{}i+\allowbreak{}1, where L\_\allowbreak{}i is now a positive power of q. Thus r\_\allowbreak{}i+\allowbreak{}s\_\allowbreak{}i=\allowbreak{}L\_\allowbreak{}i-1\textgreater{}\allowbreak{}=\allowbreak{}0 and r\_\allowbreak{}i\textasteriskcentered{}\allowbreak{}s\_\allowbreak{}i\textgreater{}\allowbreak{}=\allowbreak{}0 imply both are nonnegative: one negative would force the other nonpositive and make the sum negative. This recovery occurs after power semantics and cannot be used to justify the initial t orientation. Our initial t orientation did not use any remainder bound.



The usual bounded remainder equations then restore all quotient signs: a positive numerator, 0\textless{}\allowbreak{}=\allowbreak{}r\textless{}\allowbreak{}L, and positive divisors exclude a negative high or range quotient. The endpoint quotients are nonnegative because x\textgreater{}\allowbreak{}=\allowbreak{}0, y\textgreater{}\allowbreak{}F, and the two endpoint rows. The front is therefore genuinely natural wherever its inherited arithmetic proof requires it.



\subsection{Self-sized radix and the unchanged factorial fixed point}



From the preceding raw ceiling bootstrap, V0\textless{}\allowbreak{}=\allowbreak{}BF\textless{}\allowbreak{}B\textasciicircum{}\allowbreak{}2. Since K\textgreater{}\allowbreak{}=\allowbreak{}2 and H0\textgreater{}\allowbreak{}=\allowbreak{}1,



\begin{verbatim}
q=(B H0)^K>=B^K>=B^2>V0.
\end{verbatim}

For positive rho and t, every binomial coefficient binom(rho+\allowbreak{}t,j) with 0\textless{}\allowbreak{}=\allowbreak{}j\textless{}\allowbreak{}=\allowbreak{}t is at most (rho+\allowbreak{}1)\textasciicircum{}\allowbreak{}t; this is the inherited Pascal/\allowbreak{}product bound, which does not require rho\textgreater{}\allowbreak{}=\allowbreak{}t. Therefore the single digit is exact and B=\allowbreak{}binom(F\textasciicircum{}\allowbreak{}2+\allowbreak{}t,t). The round15 square-output fixed-point proof, valid for t\textgreater{}\allowbreak{}=\allowbreak{}16, then gives F=\allowbreak{}t!. This proof uses F\textgreater{}\allowbreak{}t\textasciicircum{}\allowbreak{}2, obtained from 2\textasciicircum{}\allowbreak{}t/\allowbreak{}(t+\allowbreak{}1)\textgreater{}\allowbreak{}t\textasciicircum{}\allowbreak{}2, to exclude any small additional fixed point.



The same B is the product of the pairwise coprime d\_\allowbreak{}i=\allowbreak{}1+\allowbreak{}F\textasciicircum{}\allowbreak{}2/\allowbreak{}i, 1\textless{}\allowbreak{}=\allowbreak{}i\textless{}\allowbreak{}=\allowbreak{}t. Thus N=\allowbreak{}B\textasciicircum{}\allowbreak{}K is still the amplified CRT modulus. The new q is N\textasteriskcentered{}\allowbreak{}H0\textasciicircum{}\allowbreak{}K, so N divides q.



Every independent representative a\_\allowbreak{}i is strictly below H0. For an ordinary lower range index 1\textless{}\allowbreak{}=\allowbreak{}k\textless{}\allowbreak{}R0 and j\textless{}\allowbreak{}=\allowbreak{}k,



\begin{verbatim}
binom(a_i,j)<=a_i^k<H0^K<=q,
\end{verbatim}

with the coefficient zero if j\textgreater{}\allowbreak{}a\_\allowbreak{}i. No aggregate representative bound is needed. For the derived label Jlabel the inherited ledger gives Jlabel\textless{}\allowbreak{}=\allowbreak{}(n+\allowbreak{}1)\textasciicircum{}\allowbreak{}ell, where n is the largest of its four syntax parameters and ell=\allowbreak{}4194304. Since H0 exceeds every positive independent representative, n+\allowbreak{}1\textless{}\allowbreak{}=\allowbreak{}H0 and Jlabel\textless{}\allowbreak{}=\allowbreak{}H0\textasciicircum{}\allowbreak{}ell. The retained padding t\textgreater{}\allowbreak{}=\allowbreak{}16 and F=\allowbreak{}t! give ell(t+\allowbreak{}1)\textless{}\allowbreak{}2F\textless{}\allowbreak{}K, hence every needed derived-label coefficient is below H0\textasciicircum{}\allowbreak{}K\textless{}\allowbreak{}=\allowbreak{}q. The earlier stronger coefficient estimate using 3ell also suffices.



The local Horn majorant still gives L\textless{}\allowbreak{}R0\textasciicircum{}\allowbreak{}8388608\textless{}\allowbreak{}2\textasciicircum{}\allowbreak{}R0\textless{}\allowbreak{}2\textasciicircum{}\allowbreak{}K. All prime-power range-decoding and multiplicity arguments therefore apply exactly as in the inherited proof. In particular they do not assume that a supplied rule multiset is acyclic, and the truth label still has unit weight.



\subsection{Cofinal completeness at the new beta modulus}



Start with a finite C222 refutation, not with an assumed solution or consistency assumption. Expand its DAG into a finite tree. Choose h0 large enough that t=\allowbreak{}16+\allowbreak{}h0\textasciicircum{}\allowbreak{}2 dominates its step count, every local parameter and every full constructor label sum, including the fixed truth-padding tuple. Insert truth steps to reach exactly t rows. This restriction is complete because every sufficiently large horizon is permitted, not merely because an unspecified admissible set is unbounded.



Set F=\allowbreak{}t!, b=\allowbreak{}F\textasciicircum{}\allowbreak{}2+\allowbreak{}1 and V0=\allowbreak{}b\textasciicircum{}\allowbreak{}t. Each pending multiset has at most t+\allowbreak{}1 members, including multiplicities, and every weight is at most 1+\allowbreak{}tF\textless{}\allowbreak{}=\allowbreak{}F\textasciicircum{}\allowbreak{}2+\allowbreak{}1=\allowbreak{}b because F\textgreater{}\allowbreak{}=\allowbreak{}t. Hence each accumulator is at most b\textasciicircum{}\allowbreak{}(t+\allowbreak{}1). Since t\textgreater{}\allowbreak{}=\allowbreak{}16 and b\textgreater{}\allowbreak{}=\allowbreak{}2,



\begin{verbatim}
y=V0(V0-1)=b^t(b^t-1)>b^(t+1).
\end{verbatim}

Moreover V0=\allowbreak{}1 modulo F, so F divides y, and y\textgreater{}\allowbreak{}=\allowbreak{}F+\allowbreak{}1. These are precisely the inherited first beta-CRT requirements. The moduli 1+\allowbreak{}(i+\allowbreak{}1)y, 0\textless{}\allowbreak{}=\allowbreak{}i\textless{}\allowbreak{}=\allowbreak{}t, are pairwise coprime: a common prime divides i-j, hence is at most t and divides F and y, contradicting its divisibility of either modulus. Choose the CRT code x and its nonnegative endpoint quotients.



Then fix R0,K,B,N and choose the nine independent positive representatives by the original CRT at precision d\_\allowbreak{}i\textasciicircum{}\allowbreak{}(2K). No aggregate cap is imposed on them. Only afterwards compute H0 and q=\allowbreak{}(B H0)\textasciicircum{}\allowbreak{}K. The bounds above validate all required digits uniformly in these chosen representatives. This dependence order has no fixed-point or future-modulus circularity.



\subsection{Complete private Pell witnesses}



For these genuine front values form the exact square frame W. Set A=\allowbreak{}X\_\allowbreak{}W(W+\allowbreak{}1). The recurrence Y\_\allowbreak{}(n+\allowbreak{}2)=\allowbreak{}2(W+\allowbreak{}1)Y\_\allowbreak{}(n+\allowbreak{}1)-Y\_\allowbreak{}n gives Y\_\allowbreak{}n(W+\allowbreak{}1)=\allowbreak{}n modulo 2W, so Y\_\allowbreak{}W(W+\allowbreak{}1)/\allowbreak{}W is a positive odd integer. Set z=\allowbreak{}(Y\_\allowbreak{}W(W+\allowbreak{}1)/\allowbreak{}W-1)/\allowbreak{}2. This supplies the odd frame exactly.



Choose each main pair canonically at k\_\allowbreak{}i, and choose the anchor U=\allowbreak{}X\_\allowbreak{}W(A), C=\allowbreak{}Y\_\allowbreak{}W(A). Let h=\allowbreak{}max\_\allowbreak{}i k\_\allowbreak{}i. The square frame gives W\textgreater{}\allowbreak{}=\allowbreak{}2h+\allowbreak{}1. Pell doubling and growth imply



\begin{verbatim}
Y_(2h)=2X_hY_h>2Y_h^2,
Y_(2h+1)>(2A-1)Y_(2h)>2(2A-1)Y_h^2.
\end{verbatim}

Thus C-sum\_\allowbreak{}i Y\_\allowbreak{}(k\_\allowbreak{}i)\textasciicircum{}\allowbreak{}2\textgreater{}\allowbreak{}(4A-2-m)Y\_\allowbreak{}h\textasciicircum{}\allowbreak{}2\textgreater{}\allowbreak{}W. Consequently the slack in (1) is nonnegative. Its ternary witnesses exist by the standard representation n=\allowbreak{}u\textasciicircum{}\allowbreak{}2+\allowbreak{}v\textasciicircum{}\allowbreak{}2+\allowbreak{}w(w+\allowbreak{}1) of every nonnegative integer.



Choose a positive Pell unit power at A that is one modulo 8C\textasciicircum{}\allowbreak{}2. Its coordinates give H odd, J\textgreater{}\allowbreak{}0 divisible by C\textasciicircum{}\allowbreak{}2 and gcd(H,8J)=\allowbreak{}1. Set e=\allowbreak{}J/\allowbreak{}C\textasciicircum{}\allowbreak{}2-1\textgreater{}\allowbreak{}=\allowbreak{}0. CRT now chooses a sufficiently large positive D with



\begin{verbatim}
D=A modulo H,       D=1+4J modulo 8J.
\end{verbatim}

Hence D=\allowbreak{}1+\allowbreak{}4J(2d+\allowbreak{}1) for an integer d. Choose the comparison pair at index W; polynomial base transfer gives its two required congruences. Fill the three signed quotient witnesses. This constructs the full single anchor, including the strengthened odd multiplier condition.



For the pooled row, each Y\_\allowbreak{}(k\_\allowbreak{}i)-k\_\allowbreak{}i is divisible by A-1, so its sum of squares is divisible by A-1. The sparse quotients exist individually. Fill every genuine power congruence by the ordinary recurrence identity and every retained sign guard by the ternary representation. This gives a complete zero of the new source. Conversely, the prior soundness argument gives a genuine front refutation; one may reconstruct old natural Pell witnesses afresh. No claim identifies the old and new private witnesses pointwise.



\subsection{Exact compilation and verification boundary}



\texttt{build\_\allowbreak{}cap.py} records the full square-frame atom list, removed private constraints, complete new anchor, pooled or individual residues, actual sign guards, front changes, and source hashes. Every guard G\textgreater{}\allowbreak{}=\allowbreak{}0 is the quadratic row G=\allowbreak{}a\textasciicircum{}\allowbreak{}2+\allowbreak{}b\textasciicircum{}\allowbreak{}2+\allowbreak{}c(c+\allowbreak{}1). Products of interval coordinates and the linear-in-anchor sum of squares are themselves quadratic.



The inherited exact Horn factoring does not alter its external interface. Monic integral projection then removes only equations z-Q=\allowbreak{}0 with integral quadratic Q and substitutions that leave every residual quadratic. Its certificate restores every deleted variable; unconstrained variables may be restored as zero. Integral translations are bijections old=\allowbreak{}new+\allowbreak{}shift. The separately stored collected polynomial is the exact sum of squares, not a numerical sample or a symbolic exponent shorthand.



For reproduction from the workspace root:



\begin{verbatim}
python diophantine_zfc_round16/direct/build_cap.py diophantine_zfc_round16/direct/portable_input diophantine_zfc_round16/direct/rebuilt --strong-cap --odd-frame --odd-comparison --ys-only --cofinal --self-radix
python diophantine_zfc_round16/direct/center_integer.py diophantine_zfc_round16/direct/rebuilt diophantine_zfc_round16/direct/rebuilt_integer324
\end{verbatim}

Add --sparse for the 338-variable companion. The separate conservative checkpoints retain stronger sign conditions or omit cofinal restrictions. Their counts must not be confused with the released two endpoints. Literal checks certify syntax, complete expansion and exact algebraic maps. The unbounded Pell, factorial, CRT, Horn-flow and calibration arguments above and in the identified dependencies supply the semantic theorem.

% Complete unabridged proof from scalar/PROOF.md
\section{Parity normalized product equations in eleven unrestricted integers}



\subsection{Result and scope}



This component gives two complete ordinary integer polynomials, each with eleven unrestricted integer unknowns and the six unchanged, fixed source34 coefficient ports. The native-index form has \textbf{109 binary operations and exact degree 966}. Replacing its first-index modulus by \texttt{4q\textasciicircum{}\allowbreak{}4} gives \textbf{110 operations and exact degree 882}. These are operation/\allowbreak{}degree tradeoffs against the independently developed lower-degree three-root forms.



The arithmetic mechanism is a product of three pairwise coprime integers. A congruence modulo four recovers their signs after square extraction. There is no radical norm in the displayed equation. The strong Pell ordinate is saturated with the square root of an already coprime product; the auxiliary coordinate is rescaled compatibly. A weak scalar carrier and the supported-modulus interpolation theorem supply the reverse direction.



The claimed consistency equivalence uses the same explicitly retained ordinary source34/\allowbreak{}native/\allowbreak{}compiler/\allowbreak{}refutation calibration as the predecessor. The new result is the complete ordinary arithmetic reduction below. It is not a new Lean/\allowbreak{}Rocq endpoint, a machine-checked object-ZFC derivation, a proof of consistency, or a global simplicity minimum.



\subsection{1. Complete polynomial}



All eleven independent variables



\begin{verbatim}
j, C, z, h, r, v, w, T, c, n, s
\end{verbatim}

range over \textbf{all integers}. The fixed ports are



\begin{verbatim}
beta, K0, MC, MF, cone_A, cone_B.
\end{verbatim}

They are the entire coherent source34 tuple, including the already-shifted MF, with \texttt{beta=\allowbreak{}2\textasciicircum{}\allowbreak{}1064921-1}, \texttt{cone\_\allowbreak{}A=\allowbreak{}2\textasciicircum{}\allowbreak{}1113}, and \texttt{cone\_\allowbreak{}B=\allowbreak{}2\textasciicircum{}\allowbreak{}666104}. The supplied compressed bytes and reconstruction recipe specify all six numerals. They are not quantified parameters.



Every following symbol is a polynomial abbreviation:



\begin{verbatim}
J=6j+1, q=beta J+1, W=12w+4,
X=qW, Y=(2v+1)q^3, E=XY,
a=(X+1)Y, A=a+2, H=4a+3, Delta=a^2+H=A^2-1,
F=(K0+W)C-z(q-1),
R=[q(q-F)-C](q^2-1)+(MC+q MF)J.
\end{verbatim}

Choose one of two explicit moduli:



\begin{verbatim}
M_index=XY              (109-operation form),
M_index=4q^4            (110-operation form).
\end{verbatim}

Continue with



\begin{verbatim}
k=R+1+h M_index, eta=c-Yk,
Derr=(cone_A F-q)^2+(cone_B C-q)^2,
Budget=k^2(q^2-Derr)-4q^2 eta^2,
D=Budget-4,
ell=XY^2, Dt=1+ell(ell+1)k^2,
Dg=1+D Dt s^2, P=Dt Dg,
Dm=1+Delta c^2,
U=r c^2 A Dm,
Df=1+Delta U^2 P,
Q=Delta^2 U^2 P, delta=Q-1,
b=c-R+4UPT, V=Df b-c, Cy=QV^2-1,
K=X(3c-2X)+2,
S=srH, zc=HT+K,
Ac=(zc^2+3H^2)(zc^3-19H^3), Z=Sn-Ac.
\end{verbatim}

The single equation is



\begin{verbatim}
delta Z^2 - S^2 P Df Cy = 0.                         (1)
\end{verbatim}

This is a division-free ordinary polynomial with integer coefficients. In the literal graph Dg is not materialized; the identical identity



\begin{verbatim}
P=Dt+D(Dt s)^2
\end{verbatim}

uses no forced factor two in its guard scale. In the short-index form, the existing register \texttt{4q\textasciicircum{}\allowbreak{}2} from the budget is also used to compute \texttt{(4q\textasciicircum{}\allowbreak{}2)q\textasciicircum{}\allowbreak{}2=\allowbreak{}4q\textasciicircum{}\allowbreak{}4}, saving one more multiplication without changing the polynomial. The shared product \texttt{UP=\allowbreak{}U\textasteriskcentered{}\allowbreak{}P} is reused in \texttt{strong=\allowbreak{}(Delta U)\textasteriskcentered{}\allowbreak{}UP} and \texttt{b=\allowbreak{}c-R+\allowbreak{}4UPT}.



\subsection{2. Unconditional nonzero facts}



The actual affine restrictions and fixed tuple give, for every integer assignment,



\begin{verbatim}
q even, q=2 mod3, |q|>=beta-1,
X=0 mod8, X=2 mod3, |X|>=4|q|,
Y=0 mod8, |Y|>=|q|^3,
|a|>2, A!=0, H odd and nonzero, Delta>1.
\end{verbatim}

In particular A is even and Dm is a positive integer. For every integer ell, ell(ell+\allowbreak{}1)\textgreater{}\allowbreak{}=\allowbreak{}0, so Dt\textgreater{}\allowbreak{}=\allowbreak{}1. Moreover ell is a multiple of eight, hence Dt=\allowbreak{}1 modulo eight. The authentic cone coefficients are divisible by sixteen. Writing Ucone=\allowbreak{}cone\_\allowbreak{}A F and Vcone=\allowbreak{}cone\_\allowbreak{}B C, expansion gives



\begin{verbatim}
Budget+(qk)^2
  = k^2[2q(Ucone+Vcone)-Ucone^2-Vcone^2]-4q^2 eta^2
  = 0 modulo sixteen.
\end{verbatim}

Here q is even, so every displayed summand is divisible by sixteen. An even square is 0 or 4 modulo sixteen. Consequently



\begin{verbatim}
Budget=0 or12 modulo16, D=12 or8 modulo16.          (2a)
\end{verbatim}

In particular D is divisible by four and is never an integer square. This is the new parity lemma: it replaces both the former scalar factor three in D and the forced factor four in Dg. Since Dg=\allowbreak{}1 modulo four, it is nonzero and



\begin{verbatim}
P=1 mod4, P!=0.                                    (2)
\end{verbatim}

Because Delta\textgreater{}\allowbreak{}1, neither \texttt{Df=\allowbreak{}1+\allowbreak{}Delta U\textasciicircum{}\allowbreak{}2 P} nor \texttt{delta=\allowbreak{}Delta\textasciicircum{}\allowbreak{}2 U\textasciicircum{}\allowbreak{}2 P-1} can vanish. Also \texttt{Cy=\allowbreak{}QV\textasciicircum{}\allowbreak{}2-1} cannot vanish: an integer identity QV\textasciicircum{}\allowbreak{}2=\allowbreak{}1 would force Q=\allowbreak{}1, already excluded. Thus



\begin{verbatim}
Da=(Q-1)(QV^2-1)=delta Cy
\end{verbatim}

is nonzero on every integer assignment. No sign or Pell interpretation has been used for these assertions.



If S=\allowbreak{}0, equation (1) reduces to \texttt{delta Ac\textasciicircum{}\allowbreak{}2=\allowbreak{}0}. This is impossible. Indeed H!=\allowbreak{}0 and



\begin{verbatim}
Ac=H^5 pi(T+K/H), pi(x)=(x^2+3)(x^3-19).
\end{verbatim}

The first factor has no real zero; the second has no rational zero. Therefore every zero of (1) has S!=\allowbreak{}0, and hence s,r are nonzero. The value c has not yet been shown nonzero.



\subsection{3. Coprimality and sign recovery, including c=\allowbreak{}0}



All gcds here mean gcds of absolute integer values. Directly,



\begin{verbatim}
gcd(P,Df)=1,
Q=0, delta=-1, Cy=-1, Da=1 modulo P,
\end{verbatim}

so gcd(P,Da)=\allowbreak{}1. Modulo Df,



\begin{verbatim}
Q=-Delta, delta=-A^2, V=-c, Cy=-Dm,
Da=A^2 Dm.
\end{verbatim}

Since \texttt{A Dm | U}, the integer Df is 1 modulo both A and Dm. Consequently



\begin{verbatim}
gcd(Df,Da)=1.                                      (3)
\end{verbatim}

These identities remain valid for negative P or Df and for U=\allowbreak{}0.



Multiplying (1) by delta and dividing by the nonzero rational square S\textasciicircum{}\allowbreak{}2 gives



\begin{verbatim}
(delta Z/S)^2 = P Df Da.                           (4)
\end{verbatim}

The right side is a nonzero integer. A rational number whose square is an integer is an integer, so it is an integer square. By (3), the absolute value of each of P,Df,Da is separately an integer square. By (2), P cannot be the negative of an odd square, which would be 3 modulo four. Thus P\textgreater{}\allowbreak{}0. Now Df=\allowbreak{}1+\allowbreak{}Delta U\textasciicircum{}\allowbreak{}2 P\textgreater{}\allowbreak{}0, and (4) forces Da\textgreater{}\allowbreak{}0. We obtain positive integers



\begin{verbatim}
p0^2=P, f^2=Df, e0^2=Da.                           (5)
\end{verbatim}

There is a second coprime split:



\begin{verbatim}
gcd(Dt,Dg)=gcd(Dt,1+D Dt s^2)=1.
\end{verbatim}

Dt is positive and P=\allowbreak{}DtDg\textgreater{}\allowbreak{}0, so Dg\textgreater{}\allowbreak{}0. Therefore



\begin{verbatim}
Dt=tau^2, Dg=g^2, p0=tau g
\end{verbatim}

for positive integers tau,g. Since s!=\allowbreak{}0, Budget\textless{}\allowbreak{}=\allowbreak{}3 would imply D\textless{}\allowbreak{}=\allowbreak{}-1 and Dg\textless{}\allowbreak{}=\allowbreak{}1-Dt\textless{}\allowbreak{}=\allowbreak{}0. Dg is nonzero, so this contradicts Dg\textgreater{}\allowbreak{}0. Hence Budget\textgreater{}\allowbreak{}=\allowbreak{}4; (2a) sharpens this to Budget\textgreater{}\allowbreak{}=\allowbreak{}12. In particular Budget\textgreater{}\allowbreak{}0.



The centered-budget argument is now valid. It supplies k!=\allowbreak{}0, Derr\textless{}\allowbreak{}q\textasciicircum{}\allowbreak{}2, the common nonzero sign of q,F,C,q-F-C, and



\begin{verbatim}
|c/k-Y|<1/2, sign(c)=sign(kY).
\end{verbatim}

It also excludes c=\allowbreak{}0 directly: at c=\allowbreak{}0,



\begin{verbatim}
Budget=k^2[q^2-Derr-4q^2Y^2]<=0,
\end{verbatim}

since |Y|\textgreater{}\allowbreak{}=\allowbreak{}1. Thus c!=\allowbreak{}0 and U!=\allowbreak{}0. We have now recovered positivity before invoking any Pell classification or dividing by c.



\subsection{4. Restore the complete strong and auxiliary source}



Define only in the proof



\begin{verbatim}
I=U p0, T_old=p0 T.
\end{verbatim}

They are integers, \texttt{I!=\allowbreak{}0}, \texttt{c\textasciicircum{}\allowbreak{}2|I}, and \texttt{|I|\textgreater{}\allowbreak{}=\allowbreak{}c\textasciicircum{}\allowbreak{}2}. The exact restored identities are



\begin{verbatim}
f^2=1+Delta I^2,
Q=(Delta I)^2>1,
b=c-R+4I T_old,
V=f^2(c-R+4I T_old)-c.
\end{verbatim}

The rational number \texttt{y=\allowbreak{}e0/\allowbreak{}delta} satisfies



\begin{verbatim}
QV^2-(Q-1)y^2=1.                                   (6)
\end{verbatim}

It is integral. Here is the adjacent-coefficient descent in the form needed here. If y=\allowbreak{}u/\allowbreak{}d is in lowest terms, d\textasciicircum{}\allowbreak{}2 divides Q-1. The map



\begin{verbatim}
V'=(2Q-1)V-2(Q-1)y,
y'=(2Q-1)y-2QV
\end{verbatim}

preserves (6), keeps V' integral, and preserves the reduced denominator d. After reflecting signs, V\textgreater{}\allowbreak{}1 and y\textgreater{}\allowbreak{}0 give V\textless{}\allowbreak{}y\textless{}\allowbreak{}QV/\allowbreak{}(Q-1), hence -V\textless{}\allowbreak{}V'\textless{}\allowbreak{}V. The norm prevents V'=\allowbreak{}0. Descent reaches |V|=\allowbreak{}1, where |y|=\allowbreak{}1. Thus d=\allowbreak{}1. This proves the assertion without assuming divisibility by delta from the fact that Da is a square.



Canceling nonzero delta in (1) after (5)--(6) gives



\begin{verbatim}
Z^2=S^2 p0^2 f^2 y^2.
\end{verbatim}

Thus Z/\allowbreak{}S is an integer, and S|Ac. Because H|S and Ac=\allowbreak{}K\textasciicircum{}\allowbreak{}5 modulo H,



\begin{verbatim}
H | K^5.                                          (7)
\end{verbatim}

The congruences needed for signed rank recovery hold exactly:



\begin{verbatim}
V=-c modulo f,
V+R=Delta I^2(c-R)+4I Df T_old=0 modulo c.
\end{verbatim}

No assumption that I/\allowbreak{}c\textasciicircum{}\allowbreak{}2 is odd has been introduced. The retained signed-folding theorem explicitly permits an arbitrary nonzero integer strong multiplier.



\subsection{5. Smaller first-index modulus}



For the native-index variant the retained centered-rank proof applies without any change to its first-index congruence. For the short variant we verify the only changed input to that argument.



The packing bounds supplied by Budget\textgreater{}\allowbreak{}0 give



\begin{verbatim}
0<R<|q|^4-2.
\end{verbatim}

Set M=\allowbreak{}4q\textasciicircum{}\allowbreak{}4. Then M\textgreater{}\allowbreak{}3R+\allowbreak{}1. The identity



\begin{verbatim}
2XY^2=2(12w+4)(2v+1)^2 q^7
\end{verbatim}

shows \texttt{M | 2XY\textasciicircum{}\allowbreak{}2}. Write xi=\allowbreak{}sign(X), sigma=\allowbreak{}sign(k), and \texttt{PellBase=\allowbreak{}|2XY\textasciicircum{}\allowbreak{}2+\allowbreak{}1|}. The latter is congruent to xi modulo M. The first Pell classification therefore yields an index n1\textgreater{}\allowbreak{}=\allowbreak{}1 and



\begin{verbatim}
|k|=2 psi_PellBase(n1),
sigma xi^(n1-1) 2n1 = R+1 modulo M.                (8)
\end{verbatim}

The unaltered strip and growth estimates give the main index p\textgreater{}\allowbreak{}=\allowbreak{}1 with \texttt{|c|=\allowbreak{}psi\_\allowbreak{}|A|(p)} and p\textgreater{}\allowbreak{}n1. If the sign in (8) is positive then n1\textgreater{}\allowbreak{}=\allowbreak{}(R+\allowbreak{}1)/\allowbreak{}2. If negative, n1\textgreater{}\allowbreak{}R because M\textgreater{}\allowbreak{}3R+\allowbreak{}1. Thus all initial growth bounds in the retained proof hold. Strong and auxiliary folding then give p=\allowbreak{}R, without using the projection (7). Since n1\textless{}\allowbreak{}p=\allowbreak{}R, the negative case is impossible. Both 2n1 and R+\allowbreak{}1 lie strictly between zero and M, so (8) gives \texttt{2n1=\allowbreak{}R+\allowbreak{}1} exactly.



All subsequent ratio, projection and sign arguments are unchanged. After their orientation conclusions, the exact first-Pell formula is



\begin{verbatim}
k=2 psi_(2XY^2+1)((R+1)/2).
\end{verbatim}

Reduction of this recurrence modulo XY gives \texttt{XY | k-R-1}. Hence \texttt{h\_\allowbreak{}old=\allowbreak{}(k-R-1)/\allowbreak{}(XY)} is an integer and restores the original source's literal h coordinate. A new h is allowed because it occurs in that source only in the definition of k.



For completeness, the canonical first-Pell value satisfies the same recurrence. Since M|2XY\textasciicircum{}\allowbreak{}2, it has \texttt{M | k-R-1}; choose the new h as the corresponding integer quotient. Thus both directions are proved. The change does not assume that XY and 4q\textasciicircum{}\allowbreak{}4 are equal or divide one another.



\subsection{6. Fifth-power projection and the finite-refutation endpoint}



The retained signed folding and rank proof, with Section 5 for the short index, now gives \texttt{R=\allowbreak{}2t+\allowbreak{}1}, \texttt{t\textgreater{}\allowbreak{}=\allowbreak{}24}, and



\begin{verbatim}
c=sigma psi_A(R), x=|X|>=4(beta-1)>=2^64,
|Y|>(x-3)^t/2-1,
|H|>x(x-3)^t>x^(t+1)/2^t.
\end{verbatim}

Let \texttt{Kquad=\allowbreak{}X(X+\allowbreak{}2ac)-1}. The identity



\begin{verbatim}
K=HXc-2Kquad
\end{verbatim}

and odd H show that (7) implies H|Kquad\textasciicircum{}\allowbreak{}5. The exact signed-Pell projector congruence is



\begin{verbatim}
2^R Kquad = Pi modulo H,
Pi=(sigma X-2^R)(2^R sigma X+1).
\end{verbatim}

Thus H|Pi\textasciicircum{}\allowbreak{}5. Each factor has absolute value less than \texttt{2\textasciicircum{}\allowbreak{}(R+\allowbreak{}1)x}, so



\begin{verbatim}
|Pi|^5<2^(20t+20)x^10<|H|.
\end{verbatim}

For the last inequality use x\textgreater{}\allowbreak{}=\allowbreak{}2\textasciicircum{}\allowbreak{}64 and



\begin{verbatim}
64(t-9)-(21t+20)=43t-596>=436>0.
\end{verbatim}

It follows that Pi=\allowbreak{}0. Its second factor is odd and nonzero, so sigma X=\allowbreak{}2\textasciicircum{}\allowbreak{}R. The retained centered orientation and half-binomial proof then restores X,c,k,Y\textgreater{}\allowbreak{}0, \texttt{0\textless{}\allowbreak{}eta/\allowbreak{}k\textless{}\allowbreak{}1/\allowbreak{}4}, the exact main-Pell projection, and every literal old signed source row after fresh private witnesses are regenerated. Section 5 restores h\_\allowbreak{}old when required.



This is the full hypothesis set of the inherited source34/\allowbreak{}native finite- refutation endpoint. Therefore a zero of either new polynomial yields a finite syntactic ZFC refutation, under that same named ordinary calibration.



\subsection{7. Completeness, with a fixed strong pair}



Start with the canonical outer witness supplied by a finite refutation. Keep j,C,z,v,w,c and the canonical first-Pell value k; choose h using the appropriate modulus. The source gives A\textgreater{}\allowbreak{}1 even, R=\allowbreak{}3 modulo four, \texttt{c=\allowbreak{}psi\_\allowbreak{}A(R)\textgreater{}\allowbreak{}0}, \texttt{c=\allowbreak{}R modulo four}, and \texttt{Dt=\allowbreak{}tau\textasciicircum{}\allowbreak{}2} with tau\textgreater{}\allowbreak{}0. It also gives the stronger centered Budget\textgreater{}\allowbreak{}=\allowbreak{}64.



The integer D=\allowbreak{}Budget-4 is positive because the canonical Budget is at least64, and it is nonsquare by (2a). The canonical Dt=\allowbreak{}tau\textasciicircum{}\allowbreak{}2 is a nonzero square, so D Dt is positive and nonsquare. Ordinary Pell existence directly gives positive integers s,g with



\begin{verbatim}
g^2=1+D Dt s^2.
\end{verbatim}

No squaring of the unit and no restriction on the parity of s is needed.



Fix these and set \texttt{p0=\allowbreak{}tau g}. Now A,Dm,p0,s,c are fixed positive integers. Choose ONE positive strong Pell pair I,f with



\begin{verbatim}
4s A Dm p0 c^2 | I,
f^2-Delta I^2=1.
\end{verbatim}

The retained finite-period construction and supported-modulus interpolation theorem supply such a pair. Define



\begin{verbatim}
r=I/(A Dm p0 c^2), U=I/p0.
\end{verbatim}

These are positive integers; r is a multiple of 4s. Keep this pair I,f fixed from now on. Every prime factor of \texttt{M=\allowbreak{}p0 s r} divides I.



The canonical projector gives H|K; put kappa=\allowbreak{}K/\allowbreak{}H. The classical intersective quintic pi(x)=\allowbreak{}(x\textasciicircum{}\allowbreak{}2+\allowbreak{}3)(x\textasciicircum{}\allowbreak{}3-19) has a root modulo every positive integer. Choose a residue a0 modulo sr with



\begin{verbatim}
pi(a0+kappa)=0 modulo sr.
\end{verbatim}

The supported-modulus interpolation theorem, applied to this same fixed I,f and modulus M, supplies an auxiliary T\_\allowbreak{}old in the prescribed class



\begin{verbatim}
T_old=p0 a0 modulo p0 s r,
\end{verbatim}

with its full integral auxiliary Pell witness y. Consequently p0|T\_\allowbreak{}old; set T=\allowbreak{}T\_\allowbreak{}old/\allowbreak{}p0. Then T=\allowbreak{}a0 modulo sr and



\begin{verbatim}
Ac/S=H^4 pi(T+kappa)/(sr)
\end{verbatim}

is an integer. The auxiliary equality is precisely the new equality, since \texttt{4I T\_\allowbreak{}old=\allowbreak{}4UPT}. Finally choose



\begin{verbatim}
n=Ac/S+p0 f y.
\end{verbatim}

This unrestricted integer makes (1) vanish. The order is outer witness, guard s/\allowbreak{}g, one strong I/\allowbreak{}f/\allowbreak{}r, intersective residue, T\_\allowbreak{}old and T/\allowbreak{}y, then n. No strong pair is changed after the residue is requested; no coprimality between H, s, r, p0, or I is assumed.



Therefore each polynomial has an integer zero exactly when the calibrated source has a finite syntactic ZFC refutation. Composing the retained ordinary calibration gives in ZFC the equivalence between the polynomial's universal nonvanishing over Z\textasciicircum{}\allowbreak{}11 and Con(ZFC).



\subsection{8. Exact degrees and coefficients}



With q0=\allowbreak{}6 beta j, X0=\allowbreak{}12q0 w, Y0=\allowbreak{}2v q0\textasciicircum{}\allowbreak{}3, a0=\allowbreak{}X0Y0, let k0 denote the top homogeneous part of k. It is



\begin{verbatim}
k0=h a0                          (native modulus),
k0=4h q0^4-q0^3 F0,
F0=12w C-q0 z                    (short modulus).
\end{verbatim}

The prefix degrees are as follows; each indicated leading part is nonzero.



\begin{verbatim}
| register | k | P | U | Df | Q | b | V | Cy | S | Ac |
|---|---:|---:|---:|---:|---:|---:|---:|---:|---:|---:|
| native | 7 | 94 | 23 | 152 | 164 | 118 | 270 | 704 | 8 | 35 |
| short | 5 | 82 | 23 | 140 | 152 | 106 | 246 | 644 | 8 | 35 |
\end{verbatim}

In both variants Z has degree 35. The first summand delta Z\textasciicircum{}\allowbreak{}2 has degree 234 or 222. The second summand has degree 966 or 882 and therefore uniquely determines the top homogeneous part.



Direct multiplication, with



\begin{verbatim}
P0=-4 q0^2 X0^4 Y0^10 k0^6 s^2,
U0=r a0^3 c^4, S0=4sr a0,
\end{verbatim}

gives the complete leading form



\begin{verbatim}
2^22 * q0^14 X0^70 Y0^112 k0^42 s^16 r^12 c^40 T^2.
\end{verbatim}

For the native form it is a single monomial. Its coefficient at



\begin{verbatim}
j^588 w^112 v^154 h^42 s^16 r^12 c^40 T^2
\end{verbatim}

is exactly \texttt{2\textasciicircum{}\allowbreak{}988 \textasteriskcentered{}\allowbreak{} 3\textasciicircum{}\allowbreak{}700 \textasteriskcentered{}\allowbreak{} beta\textasciicircum{}\allowbreak{}588\textgreater{}\allowbreak{}0}. For the short form the coefficient at



\begin{verbatim}
j^588 w^70 v^112 h^42 s^16 r^12 c^40 T^2
\end{verbatim}

is exactly \texttt{2\textasciicircum{}\allowbreak{}946 \textasteriskcentered{}\allowbreak{} 3\textasciicircum{}\allowbreak{}658 \textasteriskcentered{}\allowbreak{} beta\textasciicircum{}\allowbreak{}588\textgreater{}\allowbreak{}0}. It is selected by the maximum h exponent in k0\textasciicircum{}\allowbreak{}42. The exponent sums are respectively 966 and 882. These are polynomial degree proofs at the actual fixed beta, not degree guesses from numerical evaluation along lines.



The independent checker does not import the producer. It verifies all literal local polynomial formulas, every live gate, the complete coefficients, the exact leading monomial at every gate, the coprimality remainder identities, the auxiliary rescaling identities, and the smaller-modulus divisibility. Finite checks support the artifact and algebra; they do not replace the unbounded Pell, interpolation, or inherited consistency arguments.



\subsection{9. Read dependencies}



\begin{itemize}

\item Incoming \texttt{ZFC\_\allowbreak{}Diophantine\_\allowbreak{}Coprime\_\allowbreak{}Projection\_\allowbreak{}2026-10-24.zip}, component \texttt{integer/\allowbreak{}PROOF.md}, especially Sections 2--4; the original complete source34 coefficient tuple and its reconstruction recipe.

\item Its \texttt{integer/\allowbreak{}source/\allowbreak{}centered11-proof.md}, Sections 3, 5.2--5.4 and 6: centered signs, signed rank recovery, exact half-binomial interval, restoration to the native endpoint, and the original full-residue theorem.

\item Its \texttt{integer/\allowbreak{}source/\allowbreak{}signed-quadratic-proof.md}, Sections 2 and 5: absolute main-root folding and the signed quadratic projector congruence.

\item Incoming \texttt{ConZFC\_\allowbreak{}Powered\_\allowbreak{}Saturation\_\allowbreak{}Factorial\_\allowbreak{}Carriers\_\allowbreak{}Round14\_\allowbreak{}2026-10-24.zip}, \texttt{radicals/\allowbreak{}PROOF.md}, Sections 4--6: fifth-power saturation and supported- modulus interpolation with its explicit Pell valuation proof.

\end{itemize}



The new square-root saturation has a separately checked completeness rescaling. It is not a pointwise polynomial substitution of the old all-fused equation and does not assert a bijection of arbitrary zeros.





\subsection{10. Attribution and local closure}



The square-root saturation, three-factor coprime product, short-index modulus and witness ordering are retained from the authenticated incoming Shared Binomial Certificates lowops construction, whose complete proof is source/\allowbreak{}SHARED\_\allowbreak{}PRODUCT\_\allowbreak{}PROOF.md. They are not claimed as new in this round. The two local changes are the parity-normalized guard D=\allowbreak{}Budget-4 and the removal of the now-redundant factor four in Dg. Every affected bootstrap and completeness step has been re-proved above.



The six actual coefficient payloads, complete fixed recipe and machine/\allowbreak{}tile data are in coefficients/. The retained source34 calibration sources and ordinary signed and interpolation arguments are in inherited/. The source and coefficient bytes are pinned in SOURCE\_\allowbreak{}PINS.json. All references to other package layouts in unchanged source copies are historical; their matching retained documents are in inherited/\allowbreak{}provenance/.



The producer and independent checker are adaptations of the authenticated incoming lowops arithmetic tools. The checker does not import the producer. It verifies complete local polynomial cuts, every live gate, actual full coefficient bytes, exact leading coefficients at every DAG node, the coprimality and rescaling identities, and rejection of changed formulas. Finite checks do not replace the unbounded proof or inherited calibration.

% Complete unabridged proof from scalar/TWO_ROOT_PROOF.md
\section{Two rational square conditions with a nonzero center}



\subsection{Result and scope}



The short-index construction below is a fixed polynomial in eleven unrestricted integers of exact degree810 with112 live binary arithmetic operations. Its weak-scale companion has degree654 with117 operations. A native-index cubic-scale comparison has degree906 with111 operations but is dominated by the separate degree882/\allowbreak{}110-operation product equation and is retained only for transparent comparison.



The new mechanism removes the A Dm saturation from the incoming three-factor coprime product and instead uses a denominator-cleared two-square identity. A nonzero-center argument, rather than fixed prime-separating weights, individually recovers both rational squares. The inherited signed source calibration is unchanged. These are ordinary arithmetic proofs, not new kernel-checked or object-ZFC derivations, and no global minimum or proof of consistency is claimed.



\subsection{1. Full polynomial and its two scale choices}



The variables are j,C,z,h,r,v,w,T,c,n,s, all unrestricted integers. The six fixed numeral ports beta,K0,MC,MF,cone\_\allowbreak{}A,cone\_\allowbreak{}B are exactly the full payloads in coefficients/. In particular beta=\allowbreak{}2\textasciicircum{}\allowbreak{}1064921-1, cone\_\allowbreak{}A=\allowbreak{}2\textasciicircum{}\allowbreak{}1113, and cone\_\allowbreak{}B=\allowbreak{}2\textasciicircum{}\allowbreak{}666104. The compiler coefficient MF is already shifted. No port is quantified.



Define



\begin{verbatim}
J=6j+1, q=beta J+1, W=12w+4, X=qW,
F=(K0+W)C-z(q-1),
R=[q(q-F)-C](q^2-1)+(MC+q MF)J.
\end{verbatim}

For the degree810 equation use



\begin{verbatim}
Y=(2v+1)q^3, index_modulus=4q^4.
\end{verbatim}

For the degree906 comparison use the same Y but index\_\allowbreak{}modulus=\allowbreak{}XY. For the degree654 weak-scale equation use



\begin{verbatim}
Y=qv, index_modulus=XY.
\end{verbatim}

Continue, in all cases, with



\begin{verbatim}
E=XY, a=E+Y, H=4a+3, Delta=a^2+H=(a+2)^2-1,
k=R+1+h index_modulus, eta=c-Yk,
Derr=(cone_A F-q)^2+(cone_B C-q)^2,
Budget=k^2(q^2-Derr)-4q^2 eta^2, D=Budget-4,
ell=XY^2, Dt=1+ell(ell+1)k^2,
Dg=1+D Dt s^2, P=Dt Dg,
U=rc^2, Df=1+Delta U^2 P,
Q=Delta^2 U^2 P, delta=Q-1,
b=c-R+4UPT, V=Df b-c, Cy=QV^2-1,
K=X(3c-2X)+2.
\end{verbatim}

For the cubic-scale equations set



\begin{verbatim}
S=srH, zc=HT+K,
Ac=(zc^2+3H^2)(zc^3-19H^3).
\end{verbatim}

For the weak-scale equation set



\begin{verbatim}
M=2q^2, Dc=HM,
S=srH M^5,
zc=Dc T+MK+H(v-q^2),
Ac=(zc^2+3Dc^2)(zc^3-19Dc^3).
\end{verbatim}

Finally define, for every variant,



\begin{verbatim}
Z=S(n+V)-Ac, Aux=S^2 Cy, Guard=S^2 P Df,
Inner=delta(Z^2-Guard)-Aux,
Ftwo=Inner^2-4delta Aux Guard.                      (1)
\end{verbatim}

Equation Ftwo=\allowbreak{}0 is the complete ordinary integer polynomial. The literal graph computes P=\allowbreak{}Dt+\allowbreak{}D(Dt s)\textasciicircum{}\allowbreak{}2, U=\allowbreak{}rc\textasciicircum{}\allowbreak{}2, UP=\allowbreak{}UP, strong=\allowbreak{}(Delta U)UP, Df=\allowbreak{}strong+\allowbreak{}1 and Q=\allowbreak{}Delta strong; all are exactly the displayed polynomials. Symbols are abbreviations, not new unknowns. There are no divisions, radicals, variable exponents or unresolved ports in the graph.



\subsection{2. Unconditional facts and the new parity guard}



The fixed tuple and affine definitions give q even, q=\allowbreak{}2 modulo3, |q|\textgreater{}\allowbreak{}=\allowbreak{}beta-1, X=\allowbreak{}0 modulo8, X=\allowbreak{}2 modulo3 and |X|\textgreater{}\allowbreak{}=\allowbreak{}4|q|. H is odd and nonzero. For cubic Y, |Y|\textgreater{}\allowbreak{}=\allowbreak{}|q|\textasciicircum{}\allowbreak{}3 and Delta\textgreater{}\allowbreak{}1. For weak Y, if v=\allowbreak{}0 then Delta=\allowbreak{}3; otherwise |a|=\allowbreak{}|(X+\allowbreak{}1)qv| is much larger than4, again giving Delta\textgreater{}\allowbreak{}1. No nonzero-Y assumption is made yet in the weak case.



Since ell is an integer, ell(ell+\allowbreak{}1)\textgreater{}\allowbreak{}=\allowbreak{}0 and Dt\textgreater{}\allowbreak{}=\allowbreak{}1. Since X is a multiple of8 and Y is even in both scales, ell is a multiple of32; in particular Dt=\allowbreak{}1 modulo4. The cone coefficients are divisible by16 and q is even, so



\begin{verbatim}
Budget=-(qk)^2 modulo16,
Budget=0 or12 modulo16,
D=12 or8 modulo16.                                (2)
\end{verbatim}

For explicit verification, Budget+\allowbreak{}(qk)\textasciicircum{}\allowbreak{}2 equals



\begin{verbatim}
k^2[2q(cone_A F+cone_B C)-(cone_A F)^2-(cone_B C)^2]
    -4q^2 eta^2,
\end{verbatim}

whose terms are all divisible by16. Thus D is a multiple of4, is nonzero and is never a square. Consequently



\begin{verbatim}
Dg=1 modulo4, P=1 modulo4, P!=0.
\end{verbatim}

Since Delta\textgreater{}\allowbreak{}1, Df=\allowbreak{}1+\allowbreak{}Delta U\textasciicircum{}\allowbreak{}2 P never vanishes, Q can never be1, and Cy=\allowbreak{}QV\textasciicircum{}\allowbreak{}2-1 never vanishes. In particular delta!=\allowbreak{}0. These claims do not presuppose P\textgreater{}\allowbreak{}0 or a Pell interpretation.



The scalar quintic pi(t)=\allowbreak{}(t\textasciicircum{}\allowbreak{}2+\allowbreak{}3)(t\textasciicircum{}\allowbreak{}3-19) has no rational root. With the nonzero denominator d=\allowbreak{}H or Dc, Ac=\allowbreak{}d\textasciicircum{}\allowbreak{}5 pi(zc/\allowbreak{}d) is therefore nonzero. If S=\allowbreak{}0, (1) is delta\textasciicircum{}\allowbreak{}2 Ac\textasciicircum{}\allowbreak{}4, nonzero. Every zero has S!=\allowbreak{}0 and hence s,r!=\allowbreak{}0. This does not yet prove c!=\allowbreak{}0.



\subsection{3. The zero-center obstruction}



At Z=\allowbreak{}0 the exact polynomial identity is



\begin{verbatim}
Ftwo=S^4(Cy-delta P Df)^2.                          (3)
\end{verbatim}

Suppose Ftwo=\allowbreak{}0 and Z=\allowbreak{}0. Then Cy=\allowbreak{}delta P Df. Since Q is divisible by P, Cy=\allowbreak{}-1 modulo P. The right side is divisible by P, so P divides1. The congruence P=\allowbreak{}1 modulo4 excludes P=\allowbreak{}-1; hence P=\allowbreak{}1. Now Dt is a positive integer and P=\allowbreak{}Dt Dg=\allowbreak{}1, so Dt=\allowbreak{}Dg=\allowbreak{}1. The equality Dg=\allowbreak{}1+\allowbreak{}D Dt s\textasciicircum{}\allowbreak{}2 then forces D=\allowbreak{}0, contradicting (2). Therefore



\begin{verbatim}
Z!=0                                               (4)
\end{verbatim}

at every integer zero. No sign, square-root extraction, auxiliary denominator descent, or scale restoration has been used in proving this.



\subsection{4. Separate the two squares without a fixed weight}



Put t=\allowbreak{}Z/\allowbreak{}S, a0=\allowbreak{}Cy/\allowbreak{}delta and b0=\allowbreak{}P Df, only as rational numbers in this proof. Dividing (1) by the nonzero rational square delta\textasciicircum{}\allowbreak{}2 S\textasciicircum{}\allowbreak{}4 gives



\begin{verbatim}
(t^2-a0-b0)^2=4a0 b0.                             (5)
\end{verbatim}

By (4), t!=\allowbreak{}0. Define rational numbers



\begin{verbatim}
y0=(t^2+a0-b0)/(2t), p1=(t^2-a0+b0)/(2t).
\end{verbatim}

A direct expansion of (5) gives



\begin{verbatim}
y0^2=a0, p1^2=b0, y0+p1=t.                        (6)
\end{verbatim}

Thus the two square conditions separate algebraically. No unjustified inference that an arbitrary sum of radicals is rational is used; the nonzero-center proof supplies the essential hypothesis. Since b0 is an integer, the rational number p1 is an integer. P Df is nonzero, hence strictly positive.



Also gcd(P,Df)=\allowbreak{}1, because Df=\allowbreak{}1 modulo P. Unique prime factorization therefore makes |P| and |Df| individually integer squares. P is odd and P=\allowbreak{}1 modulo4; it cannot be the negative of an odd square. Thus P\textgreater{}\allowbreak{}0. It follows that Df=\allowbreak{}1+\allowbreak{}Delta U\textasciicircum{}\allowbreak{}2 P\textgreater{}\allowbreak{}0. Choose positive integers p0,f with



\begin{verbatim}
P=p0^2, Df=f^2.
\end{verbatim}

Since gcd(Dt,Dg)=\allowbreak{}1, Dt\textgreater{}\allowbreak{}=\allowbreak{}1 and P\textgreater{}\allowbreak{}0, their product square splits as



\begin{verbatim}
Dt=tau^2, Dg=g^2, p0=tau g
\end{verbatim}

with positive integers tau,g. If Budget\textless{}\allowbreak{}=\allowbreak{}3, then D\textless{}\allowbreak{}=\allowbreak{}-1 and s!=\allowbreak{}0 give Dg\textless{}\allowbreak{}=\allowbreak{}1-Dt\textless{}\allowbreak{}=\allowbreak{}0, contrary to Dg\textgreater{}\allowbreak{}0. Therefore Budget\textgreater{}\allowbreak{}=\allowbreak{}4; (2) strengthens this to Budget\textgreater{}\allowbreak{}=\allowbreak{}12. It follows that k!=\allowbreak{}0, Derr\textless{}\allowbreak{}q\textasciicircum{}\allowbreak{}2 and



\begin{verbatim}
|c/k-Y|<1/2, sign(c)=sign(kY) whenever cY!=0.        (7)
\end{verbatim}

The disk inequality also gives the same nonzero common signs of q,F,C,q-F-C as the retained source proof.



\subsection{5. Integer auxiliary root and carrier divisibility before scale restoration}



Set I=\allowbreak{}U p0. Then I is an integer, c\textasciicircum{}\allowbreak{}2 divides I, and



\begin{verbatim}
f^2=1+Delta I^2, Q=(Delta I)^2,
V=f^2(c-R+4I T_old)-c, T_old=p0T.                  (8)
\end{verbatim}

Equation (6) gives the rational auxiliary equation



\begin{verbatim}
QV^2-(Q-1)y0^2=1.                                 (9)
\end{verbatim}

If c=\allowbreak{}0, then U=\allowbreak{}I=\allowbreak{}Q=\allowbreak{}0 and (9) gives y0\textasciicircum{}\allowbreak{}2=\allowbreak{}1, so y0 is integral. If c!=\allowbreak{}0, then r!=\allowbreak{}0 and p0\textgreater{}\allowbreak{}0 imply I!=\allowbreak{}0, |I|\textgreater{}\allowbreak{}=\allowbreak{}c\textasciicircum{}\allowbreak{}2 and Q\textgreater{}\allowbreak{}1. Adjacent-coefficient Pell descent makes y0 integral. Namely, write y0=\allowbreak{}A/\allowbreak{}B in lowest terms. B\textasciicircum{}\allowbreak{}2 divides Q-1. The map



\begin{verbatim}
V'=(2Q-1)V-2(Q-1)y0,
y'=(2Q-1)y0-2QV
\end{verbatim}

preserves (9), keeps V' integral and preserves the least common denominator B by its integral unimodular matrix. After reflecting signs, V\textgreater{}\allowbreak{}1 and y0\textgreater{}\allowbreak{}0 imply V\textless{}\allowbreak{}y0\textless{}\allowbreak{}QV/\allowbreak{}(Q-1), hence -V\textless{}\allowbreak{}V'\textless{}\allowbreak{}V; the norm excludes V'=\allowbreak{}0. Descent reaches |V|=\allowbreak{}|y0|=\allowbreak{}1 and therefore B=\allowbreak{}1.



In both c cases, (6) now gives t=\allowbreak{}y0+\allowbreak{}p1 in Z. As Z=\allowbreak{}S(n+\allowbreak{}V)-Ac,



\begin{verbatim}
S divides Ac.                                     (10)
\end{verbatim}

This is proved before assuming the cubic scale, c!=\allowbreak{}0, or Q\textgreater{}\allowbreak{}1 in the weak-scale case. It is the critical ordering point for that variant.



For cubic Y, reducing (10) modulo H gives H|K\textasciicircum{}\allowbreak{}5. Inequality (7) and |Y|\textgreater{}\allowbreak{}=\allowbreak{}|q|\textasciicircum{}\allowbreak{}3 exclude c=\allowbreak{}0.



For weak Y, H=\allowbreak{}3 modulo q and H is odd. Since q=\allowbreak{}2 modulo3,



\begin{verbatim}
gcd(H,2q^2)=1.
\end{verbatim}

Write M=\allowbreak{}2q\textasciicircum{}\allowbreak{}2 and Dc=\allowbreak{}HM. Since M\textasciicircum{}\allowbreak{}5 divides S, (10) implies M\textasciicircum{}\allowbreak{}5|Ac. If p\textasciicircum{}\allowbreak{}e exactly divides M and v\_\allowbreak{}p(zc)=\allowbreak{}b\textless{}\allowbreak{}e, expansion



\begin{verbatim}
Ac=zc^5+3Dc^2 zc^3-19Dc^3 zc^2-57Dc^5
\end{verbatim}

shows that the first term has strictly smaller valuation5b than each other term. This contradicts p\textasciicircum{}\allowbreak{}(5e)|Ac. Thus M|zc. Reduction modulo M and gcd(H,M)=\allowbreak{}1 give



\begin{verbatim}
M | v-q^2, v=q^2(2v_old+1),
Y=(2v_old+1)q^3.                                  (11)
\end{verbatim}

The old nonzero cubic scale is restored. Inequality (7) now excludes c=\allowbreak{}0. The other component H|Ac gives H|(MK)\textasciicircum{}\allowbreak{}5, hence H|K\textasciicircum{}\allowbreak{}5 by coprimality. No unsound inference H|K from H|K\textasciicircum{}\allowbreak{}5 is made.



\subsection{6. Signed rank and native-source restoration}



After Section5, every variant has exactly the required signed kernel data:



\begin{itemize}

\item the actual cubic scale, Delta=\allowbreak{}(a+\allowbreak{}2)\textasciicircum{}\allowbreak{}2-1\textgreater{}\allowbreak{}1 and nonzero c,I;

\item c\textasciicircum{}\allowbreak{}2|I, f\textasciicircum{}\allowbreak{}2-Delta I\textasciicircum{}\allowbreak{}2=\allowbreak{}1, and |I|\textgreater{}\allowbreak{}=\allowbreak{}c\textasciicircum{}\allowbreak{}2;

\item the auxiliary Pell equation (9) with y0 integral;

\item V=\allowbreak{}-c modulo f and V=\allowbreak{}-R modulo c, by (8);

\item the first Pell square Dt=\allowbreak{}tau\textasciicircum{}\allowbreak{}2;

\item common nonzero signs of q,F,C,q-F-C, the strict centered strip (7), and the unchanged packing bounds |q|(|q|\textasciicircum{}\allowbreak{}2-1)\textless{}\allowbreak{}R\textless{}\allowbreak{}|q|\textasciicircum{}\allowbreak{}4-2;

\item the powered main projection H|K\textasciicircum{}\allowbreak{}5.

\end{itemize}



For the XY index modulus, the signed strong-index proof in inherited/\allowbreak{}centered11-proof.md Section5.2 and the absolute folding in inherited/\allowbreak{}signed-quadratic-proof.md Section2 apply. The proof explicitly allows arbitrary nonzero I/\allowbreak{}c\textasciicircum{}\allowbreak{}2; no odd multiplier is assumed. For 4q\textasciicircum{}\allowbreak{}4, the smaller-index proof in PROOF.md Section5 applies: 4q\textasciicircum{}\allowbreak{}4\textgreater{}\allowbreak{}3R+\allowbreak{}1 and 4q\textasciicircum{}\allowbreak{}4|2XY\textasciicircum{}\allowbreak{}2 restore the same first-index congruence, and after orientation the ordinary h\_\allowbreak{}old=\allowbreak{}(k-R-1)/\allowbreak{}(XY) is regenerated.



These arguments first give R=\allowbreak{}2t0+\allowbreak{}1, t0\textgreater{}\allowbreak{}=\allowbreak{}24, and



\begin{verbatim}
|H|>x(x-3)^t0, x=|X|>=2^64,
c=sign(c) psi_(a+2)(R).
\end{verbatim}

For completeness, the actual projection step uses the identity



\begin{verbatim}
2Kquad=HcX-K, Kquad=X(X+2ac)-1.
\end{verbatim}

Since H is odd, H|K\textasciicircum{}\allowbreak{}5 implies H|Kquad\textasciicircum{}\allowbreak{}5. The retained signed Pell congruence gives H|Pi\textasciicircum{}\allowbreak{}5 for



\begin{verbatim}
Pi=(sign(c)X-2^R)(2^R sign(c)X+1).
\end{verbatim}

The bounds |Pi|\textasciicircum{}\allowbreak{}5\textless{}\allowbreak{}2\textasciicircum{}\allowbreak{}(20t0+\allowbreak{}20)x\textasciicircum{}\allowbreak{}10 and |H|\textgreater{}\allowbreak{}x\textasciicircum{}\allowbreak{}(t0+\allowbreak{}1)/\allowbreak{}2\textasciicircum{}\allowbreak{}t0 imply |Pi|\textasciicircum{}\allowbreak{}5\textless{}\allowbreak{}|H|, since 64(t0-9)-(21t0+\allowbreak{}20)=\allowbreak{}43t0-596\textgreater{}\allowbreak{}=\allowbreak{}436. Thus Pi=\allowbreak{}0 and sign(c)X=\allowbreak{}2\textasciicircum{}\allowbreak{}R. The unchanged late orientation and uniform half-binomial proof restore X,c,k,Y\textgreater{}\allowbreak{}0, 0\textless{}\allowbreak{}eta/\allowbreak{}k\textless{}\allowbreak{}1/\allowbreak{}4, the exact main projector and all original native source rows after fresh private witnesses are regenerated.



This proves the soundness reduction to a finite syntactic ZFC refutation under the same named ordinary source34/\allowbreak{}compiler calibration. Its copies, provenance and notices are retained in inherited/. No new end-to-end formal verification is asserted.



\subsection{7. Completeness with no cyclic witness selection}



Start with a finite syntactic refutation and its canonical outer data. These fix q,X,Y,c,R,k,Dt=\allowbreak{}tau\textasciicircum{}\allowbreak{}2, Budget\textgreater{}\allowbreak{}=\allowbreak{}64 and the exact main projector H|K. For the short modulus choose h=\allowbreak{}(k-R-1)/\allowbreak{}(4q\textasciicircum{}\allowbreak{}4); its integrality follows from the first Pell recurrence as in PROOF.md Section5. For the weak variant set v=\allowbreak{}q\textasciicircum{}\allowbreak{}2(2v\_\allowbreak{}old+\allowbreak{}1), leaving the canonical numerical Y and all outer registers unchanged.



D=\allowbreak{}Budget-4 is positive and nonsquare by (2). Since Dt=\allowbreak{}tau\textasciicircum{}\allowbreak{}2 is a nonzero square, D Dt is positive nonsquare. Ordinary Pell existence supplies positive s,g with g\textasciicircum{}\allowbreak{}2=\allowbreak{}1+\allowbreak{}D Dt s\textasciicircum{}\allowbreak{}2. Fix these and put p0=\allowbreak{}tau g.



Choose one strong Pell pair I,f with



\begin{verbatim}
4s p0 c^2 | I, f^2-Delta I^2=1.
\end{verbatim}

The retained supported-modulus interpolation theorem permits this pair and keeps it fixed thereafter. Set r=\allowbreak{}I/\allowbreak{}(p0 c\textasciicircum{}\allowbreak{}2), so U=\allowbreak{}rc\textasciicircum{}\allowbreak{}2=\allowbreak{}I/\allowbreak{}p0. Every prime factor of p0sr divides I. No strong pair is changed after a residue is selected.



For cubic scale put kappa=\allowbreak{}K/\allowbreak{}H. Choose a0 modulo sr with pi(a0+\allowbreak{}kappa)=\allowbreak{}0 modulo sr. For weak scale use pi(a0+\allowbreak{}kappa+\allowbreak{}v\_\allowbreak{}old)=\allowbreak{}0 modulo sr. Intersectivity of pi is the retained quintic theorem. Supported-modulus interpolation for this fixed I,f supplies T\_\allowbreak{}old in the residue



\begin{verbatim}
T_old=p0 a0 modulo p0sr
\end{verbatim}

with its integral auxiliary Pell ordinate y. Consequently T=\allowbreak{}T\_\allowbreak{}old/\allowbreak{}p0 is an integer and T=\allowbreak{}a0 modulo sr. Formula (8) then is exactly the new auxiliary definition because I T\_\allowbreak{}old=\allowbreak{}UPT.



In the cubic case Ac=\allowbreak{}H\textasciicircum{}\allowbreak{}5 pi(T+\allowbreak{}kappa), giving Ac/\allowbreak{}S integral. In the weak case zc=\allowbreak{}HM(T+\allowbreak{}kappa+\allowbreak{}v\_\allowbreak{}old) and Ac=\allowbreak{}(HM)\textasciicircum{}\allowbreak{}5 pi(T+\allowbreak{}kappa+\allowbreak{}v\_\allowbreak{}old), again giving Ac/\allowbreak{}S integral. Now set



\begin{verbatim}
n=Ac/S+y+p0f-V.
\end{verbatim}

This is an integer, Z/\allowbreak{}S=\allowbreak{}y+\allowbreak{}p0f, Cy=\allowbreak{}delta y\textasciicircum{}\allowbreak{}2 and P Df=\allowbreak{}(p0f)\textasciicircum{}\allowbreak{}2. Substitution makes (1) zero. All witness choices are in the order outer data, guard s/\allowbreak{}g, one fixed I/\allowbreak{}f/\allowbreak{}r, intersective residue, T\_\allowbreak{}old and T/\allowbreak{}y, then n. The guard never depends on the subsequently chosen strong pair.



Both existence directions hold. Composing the retained ordinary calibration, ZFC proves that universal nonvanishing of each displayed fixed polynomial over Z\textasciicircum{}\allowbreak{}11 is equivalent to Con(ZFC).



\subsection{8. Exact degree and actual nonzero coefficients}



Let q0=\allowbreak{}6beta j, X0=\allowbreak{}12q0w and write Y0=\allowbreak{}2v q0\textasciicircum{}\allowbreak{}3 for cubic scale or Y0=\allowbreak{}q0v for weak scale. Let a0=\allowbreak{}X0Y0. The highest k part is h a0 for the native cubic index, 4h q0\textasciicircum{}\allowbreak{}4-q0\textasciicircum{}\allowbreak{}3F0 for the short index, and h a0-q0\textasciicircum{}\allowbreak{}3F0 for weak scale, where F0=\allowbreak{}12wC-q0z. The ties in the latter two forms are retained.



The common leading factors are



\begin{verbatim}
P0=-4q0^2 X0^4 Y0^10 k0^6 s^2,
U0=rc^2, Df0=a0^2 U0^2 P0, Q0=a0^4 U0^2 P0,
V0=4a0^2 U0^3 P0^2 T, Cy0=Q0 V0^2.
\end{verbatim}

The degree tables are



\begin{verbatim}
variant       P  Df   Q   V  Cy  S  Ac
native       94 112 124 210 544  8  35
short        82 100 112 186 484  8  35
weak         62  76  84 142 368 16  35.
\end{verbatim}

Put Wc=\allowbreak{}Sn-Ac, so Z=\allowbreak{}Wc+\allowbreak{}SV and degWc=\allowbreak{}35. The complete centering identity is



\begin{verbatim}
Inner=delta Wc^2+2delta S V Wc
            +S^2(1-V^2-delta P Df).                (12)
\end{verbatim}

In each variant the unique highest term is -S\textasciicircum{}\allowbreak{}2 V\textasciicircum{}\allowbreak{}2. Its degrees are436,388,316. The squared Inner therefore has degree872,776,632, respectively. The correction -4delta S\textasciicircum{}\allowbreak{}4 Cy P Df has degree906,810,654, respectively, strictly higher. Thus it gives the entire nonzero highest homogeneous form.



For the two cubic variants that form is



\begin{verbatim}
-2^30 q0^16 X0^50 Y0^98 k0^48 s^20 r^16 c^24 T^2.
\end{verbatim}

The native attained coefficient is



\begin{verbatim}
[j^552 w^98 v^146 h^48 s^20 r^16 c^24 T^2] Ftwo
    =-2^924 3^650 beta^552.
\end{verbatim}

The short attained coefficient is



\begin{verbatim}
[j^552 w^50 v^98 h^48 s^20 r^16 c^24 T^2] Ftwo
    =-2^876 3^602 beta^552.
\end{verbatim}

For weak scale S0=\allowbreak{}128sr a0 q0\textasciicircum{}\allowbreak{}10 and the entire top form is



\begin{verbatim}
-2^50 q0^56 X0^50 Y0^98 k0^48 s^20 r^16 c^24 T^2.
\end{verbatim}

Its maximal-h coefficient is



\begin{verbatim}
[j^300 w^98 v^146 h^48 s^20 r^16 c^24 T^2] Ftwo
    =-2^546 3^398 beta^300.
\end{verbatim}

Each monomial exponent sum is the claimed degree, and each coefficient is nonzero at the authentic fixed beta. Maximal h makes the choice unique even where k0 has two summands. These are exact actual-DAG leading-coefficient proofs, not numerical tests or generic-coefficient assumptions.



\subsection{9. Arithmetic checks and limitations}



The generators and checkers are owned adaptations of the attributed lowops DAG tools. The checkers never import their producer. They validate topology, full liveness, all eleven variables, complete six-port bindings and actual decompressed bytes, complete local polynomial cuts, the two-root and center identities, exact leading coefficients at every node with (12) at its cancellation gate, and modified-formula negative controls. Both normal and optimized Python checks pass.



The degree906 comparison is explicitly not a Pareto improvement. Degree810/\allowbreak{}112 operations and degree654/\allowbreak{}117 operations are new tradeoff points against the inspected round15 schedules. This does not assert global minimality, historical priority for Pell or norm identities, or a new formal calibration endpoint.

% Complete unabridged proof from positive/PROOF.md
\section{Three Pell square conditions compressed to one square}



\subsection{Result and scope}



There are four new fixed integer polynomials, each in ten strictly positive integer unknowns. The scaled profiles are 87 operations at exact degree 436 and 88 operations at exact degree 380. The direct profiles are 92 operations at degree 384 and 93 operations at degree 334. Operations are live binary additions, subtractions and multiplications, including fixed-constant multiplications. These improve the operation count of the 89/\allowbreak{}288 and 94/\allowbreak{}262 uncentered predecessors, not their degrees. No global minimum is asserted.



The new mechanism is simultaneous coprimality with both auxiliary-Pell factors. It replaces the four-factor two-root norm by one square equation. All coefficients are the complete, coherent compact44 BFS tuple bound in each JSON file. This is not a mixture with the public indexed tuple.



The new unbounded arithmetic argument composes with the retained ordinary native/\allowbreak{}controller/\allowbreak{}refutation calibration. Its conclusion is that ZFC proves absence of positive zeros equivalent to Con(ZFC), with the inherited finite refutation convention. It is not a new Lean theorem or an internal object-ZFC proof object, and it does not assert consistency.



\subsection{Complete polynomial definition}



The positive unknowns are J,F,alpha,i,T,w,Y,n,eta,zeta on the scaled chart; replace w by x on the direct chart. The four fixed integer ports are beta, K0,MC,MF, with beta odd, beta\textgreater{}\allowbreak{}=\allowbreak{}31, K0\textgreater{}\allowbreak{}0 and 0\textless{}\allowbreak{}MC\textless{}\allowbreak{}beta\textless{}\allowbreak{}MF\textless{}\allowbreak{}2beta. Their complete literal payloads and finite recipe are in coefficients/.



Define the following integer polynomials. No displayed intermediate is a separately quantified variable.



\begin{verbatim}
q=beta J+1, C=q-F-alpha, mu=q^2+1,
scaled: X=wq, h=q-1, B=(K0+w)C-F, P=q^3(q-1), low=(q-1)Y-q^3 B,
direct: X=q+x, h=q(q-1), B=(qK0+X)C-qF, P=q^4(q-1),
        low=q(q-1)Y-q^3 B-(q-1)X,
A=mu Y(X+1)+1, Delta=A^2-1, H=4A-5,
k=eta+zeta, c=kY+eta, Vfirst=X mu^2 Y^2,
D=1+Vfirst(Vfirst+1)k^2,
R=((q-1)(q-F)+alpha)(q^2-1)+(MC+q MF)J, e=k-R-1,
Kc=1+Delta c^2, E=i c^2 A Kc,
Fstrong=1+Delta D E^2, Q=Delta^2 D E^2, delta=Q-1,
b=c(T+alpha)-R, V=Fstrong b-c,
Dnum=Q V^2-1, G=D Fstrong, Kpos=X(3c-2X)+2,
M=P mu^2 H.
\end{verbatim}

For the 87/\allowbreak{}436 and 92/\allowbreak{}384 equations set L=\allowbreak{}V\textasciicircum{}\allowbreak{}2. For the 88/\allowbreak{}380 and 93/\allowbreak{}334 equations set L=\allowbreak{}(Delta D E)(Fstrong b). Then put



\begin{verbatim}
Z=P mu^2 [H(n-LB)+Kpos]+H[mu^2 low-Pe],
N=delta Z^2-M^2 G Dnum.
\end{verbatim}

The equation is N=\allowbreak{}0. Its complete literal circuit, without divisions, square roots, exponentiation gates or hidden witnesses, is supplied in each equations/\allowbreak{}\textasteriskcentered{}.json file. The normalized rational identity used only in proofs is



\begin{verbatim}
Z/M=n+Kpos/H+Y/q^3-(1+hL)B/h-e/mu^2-[X/q^4].             (1)
\end{verbatim}

The bracket is present only on the direct chart. Relative to the predecessor, the first, main and auxiliary kernels have the same functional form, but the strong ordinate has been additionally saturated by A Kc. There is no V or other large polynomial center in (1).



\subsection{1 Root recovery before any transport or source argument}



All positive inputs give A\textgreater{}\allowbreak{}6, c\textgreater{}\allowbreak{}=\allowbreak{}3, E\textgreater{}\allowbreak{}=\allowbreak{}c\textasciicircum{}\allowbreak{}2 A Kc, D\textgreater{}\allowbreak{}1 and delta\textgreater{}\allowbreak{}4. Fstrong\textgreater{}\allowbreak{}c+\allowbreak{}2. Because b is integral, V=\allowbreak{}Fstrong b-c is nonzero and |V|\textgreater{}\allowbreak{}1: if b=\allowbreak{}0 then V=\allowbreak{}-c; otherwise |V|\textgreater{}\allowbreak{}=\allowbreak{}Fstrong-c\textgreater{}\allowbreak{}1. Hence Dnum\textgreater{}\allowbreak{}0.



The following are polynomial congruences, valid before any square root is known to be integral:



\begin{verbatim}
Fstrong = 1 mod D, delta = -1 mod D, Dnum = -1 mod D;
delta = -A^2 mod Fstrong;
Dnum = -(1+Delta c^2) = -Kc mod Fstrong.
\end{verbatim}

The last line uses Q=\allowbreak{}Delta(Fstrong-1) and V=\allowbreak{}-c mod Fstrong. Because A and Kc divide E, Fstrong is 1 modulo both A and Kc. Thus



\begin{verbatim}
gcd(D,Fstrong delta Dnum)=1,
gcd(Fstrong,delta Dnum)=1.                              (2)
\end{verbatim}

At a zero of N, the positive integer D Fstrong delta Dnum is the rational square (delta Z/\allowbreak{}M)\textasciicircum{}\allowbreak{}2, hence is an integer square. By (2) and unique prime factorization, D and Fstrong are separately squares. Write D=\allowbreak{}tau\textasciicircum{}\allowbreak{}2 and Fstrong=\allowbreak{}f\textasciicircum{}\allowbreak{}2 with tau,f positive integers. Also delta Dnum is a square, so y=\allowbreak{}sqrt(Dnum/\allowbreak{}delta) is rational. Set I=\allowbreak{}tau E. Then



\begin{verbatim}
f^2-Delta I^2=1, Q=(Delta I)^2, c^2 divides I,
Q V^2-(Q-1)y^2=1.                                    (3)
\end{verbatim}

The retained adjacent-Pell descent makes y a positive integer. Its short proof does not assume a particular sign of V. Write y=\allowbreak{}a/\allowbreak{}m in lowest terms, so m\textasciicircum{}\allowbreak{}2 divides Q-1. For v=\allowbreak{}|V|\textgreater{}\allowbreak{}1 the transformation



\begin{verbatim}
v'=(2Q-1)v-2(Q-1)y, y'=(2Q-1)y-2Qv
\end{verbatim}

preserves Qv\textasciicircum{}\allowbreak{}2-(Q-1)y\textasciicircum{}\allowbreak{}2=\allowbreak{}1, has integral nonzero v' with |v'|\textless{}\allowbreak{}v, and preserves the reduced denominator m of y'. The inequalities needed are v\textless{}\allowbreak{}y\textless{}\allowbreak{}Qv/\allowbreak{}(Q-1). Taking absolute values and descending reaches v=\allowbreak{}1,y=\allowbreak{}1; therefore m=\allowbreak{}1. The original square equation consequently gives



\begin{verbatim}
Z/M=epsilon tau f y, epsilon in {-1,+1}.                (4)
\end{verbatim}

The single-square compression has recovered exactly the three required Pell conditions and an integral carrier. It did not assume parity, B\textgreater{}\allowbreak{}=\allowbreak{}0, the main index identity, or any source semantics.



\subsection{2 The transport majorants are sufficient}



The inherited all-index strong/\allowbreak{}auxiliary folding lemma applies to (3), with c\textless{}\allowbreak{}=\allowbreak{}I, c\textasciicircum{}\allowbreak{}2|I and V=\allowbreak{}-c mod f. It gives c=\allowbreak{}psi\_\allowbreak{}A(p) for positive odd p. The unchanged first norm gives k=\allowbreak{}2 psi\_\allowbreak{}(2Vfirst+\allowbreak{}1)(j), and j=\allowbreak{}1 is excluded as in the predecessor: it would give k=\allowbreak{}2 and c=\allowbreak{}2Y+\allowbreak{}1 strictly between the first two main ordinates. Thus k\textgreater{}\allowbreak{}=\allowbreak{}8Vfirst+\allowbreak{}4 and



\begin{verbatim}
c>8X mu^2 Y^2, Kpos>0, H>3X.
\end{verbatim}

The unconditional retained mask estimate R\textless{}\allowbreak{}c alpha/\allowbreak{}4 yields



\begin{verbatim}
b=c(T+alpha)-R>c>k,X, V>(Fstrong-1)b>b>1.
\end{verbatim}

This order does not use the sign of B. Put g=\allowbreak{}tau f and a=\allowbreak{}g y. The strong norm gives f\textless{}\allowbreak{}A I, so g\textless{}\allowbreak{}A D E. Furthermore



\begin{verbatim}
V>(Fstrong-1)c=Delta D E^2 c>4 A D E>4g,               (5)
\end{verbatim}

because Delta/\allowbreak{}A\textgreater{}\allowbreak{}A-1\textgreater{}\allowbreak{}5 and E c\textgreater{}\allowbreak{}=\allowbreak{}1. The auxiliary equation gives y\textless{}\allowbreak{}2V. Therefore a\textless{}\allowbreak{}2gV\textless{}\allowbreak{}V\textasciicircum{}\allowbreak{}2/\allowbreak{}2. Also b-e=\allowbreak{}c(T+\allowbreak{}alpha)-k+\allowbreak{}1\textgreater{}\allowbreak{}0, so



\begin{verbatim}
rho=e/mu^2<b/4,
[X/q^4]<b/4.                                        (6)
\end{verbatim}

If rho is negative the same one-sided estimate remains valid. For L=\allowbreak{}V\textasciicircum{}\allowbreak{}2, (5)-(6) show a+\allowbreak{}rho+\allowbreak{}[X/\allowbreak{}q\textasciicircum{}\allowbreak{}4]\textless{}\allowbreak{}L, since b\textless{}\allowbreak{}V\textless{}\allowbreak{}V\textasciicircum{}\allowbreak{}2.



For the lower-degree variant, SE=\allowbreak{}Delta D E\textgreater{}\allowbreak{}4g, so L=\allowbreak{}SE(Fstrong b)\textgreater{}\allowbreak{}4gV\textgreater{}\allowbreak{}2a and also L-a\textgreater{}\allowbreak{}2gV\textgreater{}\allowbreak{}b. This again gives a+\allowbreak{}rho+\allowbreak{}[X/\allowbreak{}q\textasciicircum{}\allowbreak{}4]\textless{}\allowbreak{}L. All inequalities are strict and pre-transport.



Solve (1) and (4) for the positive coordinate n:



\begin{verbatim}
n=epsilon a-Kpos/H-Y/q^3+(1+hL)B/h+rho+[X/q^4].       (7)
\end{verbatim}

If B\textless{}\allowbreak{}=\allowbreak{}-1, the transport term is at most -L-1/\allowbreak{}h. Equations (5)-(6), Kpos/\allowbreak{}H\textgreater{}\allowbreak{}0 and Y/\allowbreak{}q\textasciicircum{}\allowbreak{}3\textgreater{}\allowbreak{}0 make the right side negative. Hence B\textgreater{}\allowbreak{}=\allowbreak{}0. Since F\textgreater{}\allowbreak{}0, its definition on either chart forces C\textgreater{}\allowbreak{}0.



\subsection{3 Later source reconstruction}



This is the precise inherited dependency boundary. B\textgreater{}\allowbreak{}=\allowbreak{}0 gives 3q\textless{}\allowbreak{}R\textless{}\allowbreak{}q\textasciicircum{}\allowbreak{}4-2. The normalized strong rank and signed auxiliary-index lemmas, using c\textasciicircum{}\allowbreak{}2|I and the same congruences V=\allowbreak{}-c mod f and V=\allowbreak{}-R mod c, give p=\allowbreak{}R. Odd R gives even q and odd J. Only then does gcd(mu\textasciicircum{}\allowbreak{}2,PH)=\allowbreak{}1 isolate mu\textasciicircum{}\allowbreak{}2|e from the integral carrier (4). The first-index congruence yields j=\allowbreak{}(R+\allowbreak{}1)/\allowbreak{}2. The zero-offset ratio estimates give Y\textgreater{}\allowbreak{}X\textasciicircum{}\allowbreak{}r/\allowbreak{}3 and X\textasciicircum{}\allowbreak{}(r+\allowbreak{}1)\textless{}\allowbreak{}H, where R=\allowbreak{}2r+\allowbreak{}1, r\textgreater{}\allowbreak{}=\allowbreak{}48.



The carrier still gives H|P Kpos. The retained whole-product projection, without cancelling a nonunit through H, then gives X=\allowbreak{}2\textasciicircum{}\allowbreak{}R. The direct chart's later q|X and q\textasciicircum{}\allowbreak{}4|X steps remain in their original order. The remaining fractions give q\textasciicircum{}\allowbreak{}3|Y and the nonnegative integer t=\allowbreak{}B/\allowbreak{}h (after B=\allowbreak{}q Bt on the direct chart). The only required property of the changed weight is 1+\allowbreak{}hL=\allowbreak{}1 mod h, which holds for every integer L. The fixed-minus parity, half-binomial rounding and fresh unscaled source reconstruction are unchanged. Consequently a positive zero yields a finite refutation under the whole calibrated coefficient tuple.



Replacing a sum of roots by their product does not assert any root-index identity not already obtained from the three recovered equations. The later argument uses only integrality of the normalized carrier, not the particular integer on its right side. Full statements of the inherited interfaces and their original proofs are retained in provenance/.



\subsection{4 Fresh positive completeness}



Starting from a finite refutation, take the canonical outer and first-Pell data supplied by the inherited source theorem: X=\allowbreak{}2\textasciicircum{}\allowbreak{}R, R=\allowbreak{}3 mod 4, r\textgreater{}\allowbreak{}=\allowbreak{}48, q\textasciicircum{}\allowbreak{}3|Y, t=\allowbreak{}B/\allowbreak{}h\textgreater{}\allowbreak{}=\allowbreak{}0, c=\allowbreak{}psi\_\allowbreak{}A(R), D=\allowbreak{}tau\textasciicircum{}\allowbreak{}2, eta,zeta\textgreater{}\allowbreak{}0, and rho=\allowbreak{}(k-R-1)/\allowbreak{}mu\textasciicircum{}\allowbreak{}2\textgreater{}\allowbreak{}=\allowbreak{}1. Do not retain the former strong witness.



Choose a positive Pell index m such that



\begin{verbatim}
tau c^2 A Kc divides psi_A(m).
\end{verbatim}

It exists because the determinant-one recurrence matrix is invertible modulo this fixed positive modulus, and a positive power is the identity in the finite matrix group. Let I=\allowbreak{}psi\_\allowbreak{}A(m), f=\allowbreak{}chi\_\allowbreak{}A(m), and i=\allowbreak{}I/\allowbreak{}(tau c\textasciicircum{}\allowbreak{}2 A Kc)\textgreater{}\allowbreak{}0. This gives exactly E=\allowbreak{}i c\textasciicircum{}\allowbreak{}2 A Kc and (3).



Apply the inherited auxiliary-completeness lemma to this fresh strong pair, obtaining the same positive T and positive y,V with (3). All source congruences, including the construction of T=\allowbreak{}Tstar-alpha\textgreater{}\allowbreak{}0, are unaffected by adding a divisor to the strong period. Set the actual L to either specified polynomial and define the integer



\begin{verbatim}
n=tau f y-Kpos/H-Y/q^3+(1+hL)t+rho+[X/q^4].            (8)
\end{verbatim}

The canonical estimates give tau\textgreater{}\allowbreak{}Y\textgreater{}\allowbreak{}=\allowbreak{}Y/\allowbreak{}q\textasciicircum{}\allowbreak{}3 and f\textgreater{}\allowbreak{}Xc\textgreater{}\allowbreak{}Kpos/\allowbreak{}H. Since tau,f\textgreater{}\allowbreak{}=\allowbreak{}3 and y\textgreater{}\allowbreak{}=\allowbreak{}1, tau f y\textgreater{}\allowbreak{}=\allowbreak{}tau f\textgreater{}\allowbreak{}tau+\allowbreak{}f\textgreater{}\allowbreak{}Y/\allowbreak{}q\textasciicircum{}\allowbreak{}3+\allowbreak{}Kpos/\allowbreak{}H. All remaining terms in (8) are nonnegative, so n\textgreater{}\allowbreak{}0, even when t=\allowbreak{}0. The carrier equals tau f y, proving N=\allowbreak{}0. Every quantified coordinate is strictly positive. This is a fresh existence construction, not a claim of a positive pointwise inverse from the predecessor's zero set.



\subsection{5 Exact global degree and literal counts}



All variables have degree one and all four coefficient ports degree zero. The following degrees and leading forms follow directly from the literal definitions. Let q0=\allowbreak{}beta J, k0=\allowbreak{}eta+\allowbreak{}zeta, X0=\allowbreak{}q0 w on the scaled chart and X0=\allowbreak{}q0+\allowbreak{}x on the direct chart. Put c0=\allowbreak{}k0 Y, A0=\allowbreak{}q0\textasciicircum{}\allowbreak{}2 Y X0, D0=\allowbreak{}q0\textasciicircum{}\allowbreak{}8 X0\textasciicircum{}\allowbreak{}2 Y\textasciicircum{}\allowbreak{}4 k0\textasciicircum{}\allowbreak{}2, E0=\allowbreak{}i A0\textasciicircum{}\allowbreak{}3 c0\textasciicircum{}\allowbreak{}4. Then the leading forms of Fstrong,Q,b,V,M are respectively



\begin{verbatim}
F0=A0^2 D0 E0^2, Q0=A0^4 D0 E0^2,
b0=-q0^3(q0-F), V0=F0 b0,
M0=4q0^8 A0 (scaled), 4q0^9 A0 (direct).
\end{verbatim}

The B leading forms are B0=\allowbreak{}w(q0-F-alpha) on the scaled chart and B0=\allowbreak{}(q0 K0+\allowbreak{}X0)(q0-F-alpha)-q0 F on the direct chart. Each is nonzero, including for the actual fixed tuple; on the direct chart the coefficient of x alpha is -1.



\begin{verbatim}
register        scaled   direct
A                  5        4
Delta             10        8
c                  2        2
Kc                14       12
E                 24       21
D                 18       16
Fstrong           76       66
Q,delta           86       74
b                  4        4
V                 80       70
Dnum             246      214
G                 94       82
M                 13       13
B                  2        2
\end{verbatim}

For the low-operation variant, L0=\allowbreak{}V0\textasciicircum{}\allowbreak{}2 and deg L=\allowbreak{}160 or140. The uniquely highest part of Z is -M0 L0 B0, of degree175 or155. Consequently delta Z\textasciicircum{}\allowbreak{}2 has degree436 or384 and leading form Q0 M0\textasciicircum{}\allowbreak{}2 L0\textasciicircum{}\allowbreak{}2 B0\textasciicircum{}\allowbreak{}2. The other summand M\textasciicircum{}\allowbreak{}2 G Dnum has degree366 or322, strictly smaller. Thus N has exactly degree436 or384.



For the lower-degree variant, L0=\allowbreak{}A0\textasciicircum{}\allowbreak{}2 D0 E0 F0 b0, of degree132 or115. The same unique-dominance calculation gives deg Z=\allowbreak{}147 or130 and deg N=\allowbreak{}380 or334. These exceed366 or322 respectively. The same product Q0 M0\textasciicircum{}\allowbreak{}2 L0\textasciicircum{}\allowbreak{}2 B0\textasciicircum{}\allowbreak{}2 is the nonzero leading form. The domain is an integral polynomial ring, so these are global attained degree proofs for the actual fixed coefficients, not degrees inferred from a sampled specialization.



The literal source-row comparison is: add five saturation gates, delete the one n+\allowbreak{}V gate, and replace the thirteen-gate two-root tail by a seven-gate square tail, saving two gates overall. V\textasciicircum{}\allowbreak{}2 was already live. The lean majorant adds one multiplication. Independent topological and backward-live checks ensure every counted row and every input is used.



\subsection{6 Finite checks and their limits}



build.py reads the frozen parent JSON, never imports its source programs, and constructs the four complete graphs. verify.py independently checks the literal definitions, normalized carrier, all coefficient payload hashes, coprimality identities, exact polynomial replays and leading terms, and dependency liveness. Normal and optimized Python runs must agree. These finite checks supplement the unbounded proof; they do not establish the absence of zeros or replace the inherited logical calibration.

% Complete unabridged proof from root_math/SIGNED_POOLED_CAP.md
\section{A signed pooled Pell index certificate}



This is a new arithmetic lemma for possible integration into the direct proof-rule equation. It is not by itself a completed consistency equation. The unsigned pooled-cap idea is inherited from the authenticated incoming Shared Pell Bound 261 package. The additional contribution here is recovering both signs of every principal Pell pair through the same bounded congruence.



\subsection{Statement}



All parameters and coordinates in this note are integers. For A \textgreater{}\allowbreak{}=\allowbreak{} 2 define X\_\allowbreak{}0=\allowbreak{}1, Y\_\allowbreak{}0=\allowbreak{}0 and



\begin{verbatim}
X_(j+1)=A X_j+(A^2-1)Y_j,
Y_(j+1)=X_j+A Y_j.
\end{verbatim}

Let W\textgreater{}\allowbreak{}=\allowbreak{}4 and 1\textless{}\allowbreak{}=\allowbreak{}m\textless{}\allowbreak{}=\allowbreak{}W. Suppose A-1\textgreater{}\allowbreak{}4W\textasciicircum{}\allowbreak{}3, V=\allowbreak{}Y\_\allowbreak{}W(A), and 1\textless{}\allowbreak{}=\allowbreak{}k\_\allowbreak{}i\textless{}\allowbreak{}W for every i=\allowbreak{}1,...,m. All x\_\allowbreak{}i,y\_\allowbreak{}i,Q,u,v,w below are unrestricted integers. Impose



\begin{verbatim}
x_i^2-(A^2-1)y_i^2=1                    for every i,
V^2-1-sum_i y_i^2=u^2+v^2+w(w+1),
sum_i [(y_i-k_i)^2+(x_i-1)^2]=(A-1)Q.
\end{verbatim}

Then x\_\allowbreak{}i=\allowbreak{}X\_\allowbreak{}(k\_\allowbreak{}i)(A), y\_\allowbreak{}i=\allowbreak{}Y\_\allowbreak{}(k\_\allowbreak{}i)(A) for every i. Conversely all these canonical pairs extend to a solution of the two aggregate rows.



\subsection{Soundness}



The ternary form u\textasciicircum{}\allowbreak{}2+\allowbreak{}v\textasciicircum{}\allowbreak{}2+\allowbreak{}w(w+\allowbreak{}1) is nonnegative at every integer tuple. Thus y\_\allowbreak{}i\textasciicircum{}\allowbreak{}2\textless{}\allowbreak{}V\textasciicircum{}\allowbreak{}2, so |y\_\allowbreak{}i|\textless{}\allowbreak{}V. The complete elementary Pell classification over the integers gives a unique j\_\allowbreak{}i\textgreater{}\allowbreak{}=\allowbreak{}0 and signs sigma\_\allowbreak{}i,epsilon\_\allowbreak{}i in \{1,-1\} with x\_\allowbreak{}i=\allowbreak{}sigma\_\allowbreak{}i X\_\allowbreak{}(j\_\allowbreak{}i)(A) and y\_\allowbreak{}i=\allowbreak{}epsilon\_\allowbreak{}i Y\_\allowbreak{}(j\_\allowbreak{}i)(A). The sign epsilon\_\allowbreak{}i is immaterial when j\_\allowbreak{}i=\allowbreak{}0. Strict monotonicity of Y\_\allowbreak{}j gives 0\textless{}\allowbreak{}=\allowbreak{}j\_\allowbreak{}i\textless{}\allowbreak{}W.



Modulo A-1 the recurrence gives X\_\allowbreak{}j=\allowbreak{}1 and Y\_\allowbreak{}j=\allowbreak{}j. Hence A-1 divides



\begin{verbatim}
E=sum_i [(epsilon_i j_i-k_i)^2+(sigma_i-1)^2].
\end{verbatim}

This is a sum of nonnegative integers. Since both j\_\allowbreak{}i and k\_\allowbreak{}i belong to [0,W-1],



\begin{verbatim}
0<=E<=4m((W-1)^2+1)<=4W((W-1)^2+1)<4W^3<A-1.
\end{verbatim}

Therefore E=\allowbreak{}0. Every sigma\_\allowbreak{}i=\allowbreak{}1 and epsilon\_\allowbreak{}i j\_\allowbreak{}i=\allowbreak{}k\_\allowbreak{}i\textgreater{}\allowbreak{}0, so epsilon\_\allowbreak{}i=\allowbreak{}1 and j\_\allowbreak{}i=\allowbreak{}k\_\allowbreak{}i. No initial sign guard on x\_\allowbreak{}i,y\_\allowbreak{}i or Q was used. In particular the zero-index and negative-coordinate branches are excluded by the argument, not omitted from the classification.



\subsection{Completeness}



For canonical pairs, every summand in the pooled congruence is divisible by A-1 (indeed by its square), so its quotient Q is an integer. Every k\_\allowbreak{}i\textless{}\allowbreak{}W gives y\_\allowbreak{}i\textless{}\allowbreak{}=\allowbreak{}Y\_\allowbreak{}(W-1). The recurrence implies



\begin{verbatim}
Y_W> (2A-1)Y_(W-1).
\end{verbatim}

Since A-1\textgreater{}\allowbreak{}4W\textasciicircum{}\allowbreak{}3 and m\textless{}\allowbreak{}=\allowbreak{}W, (2A-1)\textasciicircum{}\allowbreak{}2\textgreater{}\allowbreak{}m+\allowbreak{}1. Consequently



\begin{verbatim}
V^2-1-sum_i y_i^2 >= 0.
\end{verbatim}

Every nonnegative integer G is u\textasciicircum{}\allowbreak{}2+\allowbreak{}v\textasciicircum{}\allowbreak{}2+\allowbreak{}w(w+\allowbreak{}1). To see this, Legendre's three-square theorem represents 4G+\allowbreak{}1 as a sum of three squares: its residue 1 modulo4 guarantees exactly one odd square. Write the even squares as (2u)\textasciicircum{}\allowbreak{}2,(2v)\textasciicircum{}\allowbreak{}2 and the odd square as (2w+\allowbreak{}1)\textasciicircum{}\allowbreak{}2, then divide by4. This supplies the aggregate slack witnesses.



\subsection{Polynomial and compilation boundary}



The displayed A and V are semantic parameters at this stage. An ordinary equation must still certify V=\allowbreak{}Y\_\allowbreak{}W(A), and establish A-1\textgreater{}\allowbreak{}4W\textasciicircum{}\allowbreak{}3 and all index bounds before using this lemma. In the inherited square-frame construction A\textgreater{}\allowbreak{}=\allowbreak{} (W+\allowbreak{}1)\textasciicircum{}\allowbreak{}W suffices: for W\textgreater{}\allowbreak{}=\allowbreak{}4,



\begin{verbatim}
(W+1)^W-1 >= (W+1)^4-1 >4W^3.
\end{verbatim}

The requested indices may include doubled ones provided the actual frame bounds 2k\_\allowbreak{}i\textless{}\allowbreak{}W. This must be checked from its literal atom ledger.



All displayed rows are polynomial. A quartic final sum of squares additionally requires quadratic gates for A\textasciicircum{}\allowbreak{}2 and the products in the individual Pell rows. Existing square registers can be reused in the pooled row, which is quadratic in the displayed coordinates, without expanding a large square of a long sum unnecessarily. The anchor certificate remains an explicit proof obligation. No assumption of consistency or a refutation is used.



\subsection{A stronger aggregate sharing the anchor lower bound}



Suppose in addition W\textgreater{}\allowbreak{}=\allowbreak{}2 max\_\allowbreak{}i k\_\allowbreak{}i+\allowbreak{}1. Replace the square-cap row by



\begin{verbatim}
V=W+1+sum_i y_i^2+u^2+v^2+w(w+1).
\end{verbatim}

This implies V\textgreater{}\allowbreak{}W and |y\_\allowbreak{}i|\textless{}\allowbreak{}sqrt(V)\textless{}\allowbreak{}V, so both the preceding soundness argument and the anchor's required bound V\textgreater{}\allowbreak{}=\allowbreak{}W hold at once. It removes the need for a separate anchor-lower-bound predicate.



For completeness put k=\allowbreak{}max\_\allowbreak{}i k\_\allowbreak{}i. The addition identity gives Y\_\allowbreak{}(2k)=\allowbreak{}2X\_\allowbreak{}k Y\_\allowbreak{}k\textgreater{}\allowbreak{}2Y\_\allowbreak{}k\textasciicircum{}\allowbreak{}2. Since W\textgreater{}\allowbreak{}=\allowbreak{}2k+\allowbreak{}1,



\begin{verbatim}
Y_W >= Y_(2k+1) > (2A-1)Y_(2k) > 2(2A-1)Y_k^2.
\end{verbatim}

Thus Y\_\allowbreak{}W-sum\_\allowbreak{}i y\_\allowbreak{}i\textasciicircum{}\allowbreak{}2 \textgreater{}\allowbreak{} (4A-2-m)Y\_\allowbreak{}k\textasciicircum{}\allowbreak{}2 \textgreater{}\allowbreak{} W, since m\textless{}\allowbreak{}=\allowbreak{}W, A-1\textgreater{}\allowbreak{}4W\textasciicircum{}\allowbreak{}3, and Y\_\allowbreak{}k\textgreater{}\allowbreak{}=\allowbreak{}1. Its difference from W+\allowbreak{}1 is a nonnegative integer and has the required ternary representation.



For the square frame W=\allowbreak{}2(2+\allowbreak{}sum\_\allowbreak{}a a\textasciicircum{}\allowbreak{}2), both an atom k and a doubled atom 2k meet this stronger bound: W-4|k|\textgreater{}\allowbreak{}=\allowbreak{}2((|k|-1)\textasciicircum{}\allowbreak{}2+\allowbreak{}1)\textgreater{}\allowbreak{}0. The literal compiler must include every relevant exponent atom in this sum.



\subsection{Signed index extension for bootstrap ordering}



If some requested indices have unknown sign before factorial recovery, replace 1\textless{}\allowbreak{}=\allowbreak{}k\_\allowbreak{}i\textless{}\allowbreak{}W by |k\_\allowbreak{}i|\textless{}\allowbreak{}W. The identical bounded-error proof gives x\_\allowbreak{}i=\allowbreak{}X\_\allowbreak{}(|k\_\allowbreak{}i|), and y\_\allowbreak{}i=\allowbreak{}sign(k\_\allowbreak{}i)Y\_\allowbreak{}(|k\_\allowbreak{}i|), with y\_\allowbreak{}i=\allowbreak{}0 when k\_\allowbreak{}i=\allowbreak{}0. It does not by itself prove k\_\allowbreak{}i\textgreater{}\allowbreak{}0. An application may first interpret an independently positive t-index, derive the factorial sign from its raw ceiling rows, and only then prove the remaining indices positive. This ordering must appear in the full compilation proof.



The stronger aggregate's completeness remains valid if all k\_\allowbreak{}i=\allowbreak{}0: all y\_\allowbreak{}i vanish and Y\_\allowbreak{}W\textgreater{}\allowbreak{}W+\allowbreak{}1. Otherwise the preceding proof applies with k=\allowbreak{}max\_\allowbreak{}i |k\_\allowbreak{}i|. This signed-index extension does not assert that n\textasciicircum{}\allowbreak{}k is an integer when k\textless{}\allowbreak{}0; that assertion is neither needed nor licensed.

\section{The pinned comparison and honest Pareto claims}
All quartic numbers in this table were read from complete literal artifacts and independently expanded. The new rows and their immediate parents use the same measurement convention. A dash means the construction is not presented as an ordinary integer sum of residual squares.
\begin{center}\small
\begin{tabular}{@{}lrrrr@{}}\toprule
Domain and construction&Variables&Rows&Terms&Height\\\midrule
$\N$ round 15&329&253&12210&86844780\\
$\N$ incoming shared cap&261&194&31679&85566955\\
$\N$ incoming sparse cap&292&225&6140&83933334\\
$\N$ new pooled cap&249&186&27078&85178741\\
$\N$ new sparse cap&278&215&5524&83852788\\
$\N$ incoming low height&320&255&13828&15248812\\
$\N$ incoming sparse residue&333&264&4357&209508384\\\midrule
$\Z$ round 15&442&290&4652&128446\\
$\Z$ new pooled cap&324&208&5808&128446\\
$\Z$ new sparse cap&338&222&4493&128446\\
$\Z$ incoming Weak Carriers&478&305&5484&511326\\
$\Z$ incoming squared paired guards&489&308&6905&11302118\\
$\Z$ incoming paired guards&493&310&7174&10945213\\
$\Z$ incoming cofinal witnesses&508&308&6545&996521\\
$\Z$ incoming height one&7431&---&34086&1\\\bottomrule
\end{tabular}\end{center}
The new $249$-natural equation is not uniformly better than the $278$-natural one, and neither dominates the lower-height or lower-support incoming alternatives. The $338$-integer equation dominates the round-15 direct endpoint in variables, rows and support, with equal height. The $324$-integer point has fewer variables and rows but more monomials. The height-one example is a different coefficient objective: its polynomial has only coefficients $-1,+1$, constant coefficient $1$, and individual degree at most two; $2P$, rather than $P$, has the supplied integer SOS identity. The collection also contains a $37$-integer quartic with large coefficients, so these direct-branch counts must not be presented as global fewest-variable quartics.

\subsection{The small-variable frontier}
At eleven unrestricted inputs, the selected new degree/operation pairs are $(966,109)$, $(882,110)$, $(810,112)$, and $(654,117)$. The round-15 pairs $(416,125)$, $(456,124)$, $(464,119)$ and $(488,118)$ remain relevant. The two new product points improve the incoming $(966,111)$ and $(882,112)$ at matched degree. The native two-root $(906,111)$ is dominated by $(882,110)$ and is not a selected deliverable.

At ten strictly positive inputs, the public pinned frontier includes operation/degree pairs $89/288$, $91/283$, $92/278$, $94/262$, $96/258$, and $97/254$. The complete canonical one-page display is $92/283$; the $89/288$ page is a separately preserved comparison, not the canonical display. The new $87/436$ and $88/380$ lower operation cost at higher degree. The direct-chart $92/384$ and $93/334$ artifacts are retained as complete mechanism controls, not frontier claims. No assertion transfers these positive-domain metrics to unrestricted integers.

\subsection{Review scope and provenance}
The incoming files were Cofinal Witnesses, Height One and Degree 464, Paired Guards Centered Carriers, Residue Compiler, Shared Binomial Certificates, Shared Pell Bound 261, Weak Carriers and Shared Signs, and the round-15 deliverable. The review records full exact expansion of ten incoming SOS quartics, a separate reconstruction of the height-one half-SOS polynomial, and live-operation counts for 23 integer arithmetic graphs. It separately verifies 321 source-row pullbacks, 96 constant definitions, 155 scale definitions, 572 blocks and all 7431 coordinates in the height-one compiler. These are finite algebraic audits, not independent proofs of every transitive semantic theorem.

The full review and its stated limits are in \nolinkurl{review/ROUND16_REVIEW.md}. It distinguishes proof review from count verification and explicitly identifies work assigned to the final construction audits. The report credits inherited mechanisms and records concurrent overlap; coincident metrics do not establish priority. Complete pinned URLs and byte receipts are preserved. The public snapshot is available at
\begin{center}\small
\url{https://github.com/VladimirReshetnikov/ProveIt/commit/f4d8db99e060339f3306a04d507c35bc1d0d68c5}.
\end{center}

\section{Reproduction and the evidence boundary}
The portable archive contains one final PDF, every TeX input, exact equations and coefficient payloads, complete new proofs, inherited mathematical and runtime dependencies, source pins, licenses and finite audit programs. Its top-level manifest binds every delivered file other than itself by byte size and SHA-256. Original ZIPs, old PDFs, transient build logs, caches, and unselected experimental endpoints are excluded. A historical source filename or comparison receipt is provenance, not an instruction to fetch a missing runtime dependency.

\subsection{Software and commands}
The finite replay uses Python 3.10 or later and its standard library. It makes no network request and does not require an original workspace, sibling checkout, proof assistant, or symbolic-algebra installation. All arithmetic helpers used at runtime are supplied. From the unpacked archive root:
\begin{verbatim}
python verify.py --report replay.json
python -O verify.py --report replay-optimized-driver.json
sh build_report.sh
\end{verbatim}
The verifier first authenticates its manifest, makes a fresh isolated copy, and runs builders and checkers there. Both normal and optimized interpreters are exercised by default. Explicit checks remain active under optimization. The report build uses pdfLaTeX with three passes; only the final PDF and editable sources are packaged. The included release receipt records exact command counts, byte-equality counts, hashes, and document checks.

The source34 coefficient recipe is imported through its checked finite construction function; it reconstructs the complete six-port tuple in both normal and optimized paths. Its historical standalone command-line guard is retained, not silently edited. The positive coefficient recipe is likewise reconstructed and compared to all four payloads. No port is merely a digest of an absent integer.

\subsection{What is checked independently}
Natural checks transcribe the fused factorial rows, self-sized radix, cap constraints, seeds, index additions and exact sum of squares without importing the producer. Reconstruction separately compares all emitted JSON bytes. The proof explains why the raw front is acyclic and why the horizon-twelve factorial and majorant estimates apply; sampled arithmetic is not substituted for those arguments.

The signed audit binds the original raw input, every front change, the exact square frame, one full anchor, each guard, the ordinate-only pooled or sparse congruence, Horn factoring, integral monic projections, coordinate translations, and every collected coefficient. The conservative both-sign cap is separately audited and distinguished from the released power-oriented compiler. Mathematical completeness rebuilds private Pell witnesses; no unjustified pointwise inverse is asserted.

Small-variable checks compare complete literal formulas, topology and backward liveness, all coefficient bytes, coprimality and rescaling identities, exact leading forms and intentional corruptions. Separate independent programs audit the new product and two-root or positive graphs. Finite polynomial rays are controls on transcription and the stated leading form; the unbounded attained-degree proof is given in the full text, not inferred solely from those rays.

The recorded incoming comparison audits read literal JSON data and independently reconstructed their finite measurements without executing imported programs. Their bulky input archives are omitted from this compact release, so these historical comparison receipts are not replayed by the top-level verifier. Their proof-review scope is separately stated. Passing a source hash or a finite test is not acceptance of an arithmetic theorem, and none of the checks decides consistency.

\subsection{Source closure and attribution}
The natural closure retains the exact imported executable graph, raw calculus data and ordinary source proofs. The signed closure includes its two full certificate-bearing sources, original raw input and all transformation helpers. The scalar and positive components each have a selected-file whitelist and retain their original coefficient calibration and license notices. The top-level source inventory is the complete packaging whitelist, while the manifest authenticates its actual delivered bytes.

Some unchanged historical proofs refer to their original directory layouts. Matching retained documents are supplied under the component's inherited, provenance or proofs directories; the archive's closure map gives the local locations. No old archive is required to state a final polynomial or run its finite replay. The retained calibration source files are not presented as a newly compiled formal development.

\section{Further questions}
\begin{enumerate}
\item Can the natural $K=R$ amplification and horizon-twelve proof be transferred to the signed cofinal chart without increasing the support of the literal power-oriented compiler? A correct transfer must re-establish every raw sign and carry bound before using its semantic power values.
\item Can the single anchor's comparison circuit be reduced while retaining a complete signed index proof? Eliminating its own certification would make the cap circular; finite agreement is not enough.
\item Can pooled index errors be partitioned into a small number of groups to improve the variable/support Pareto curve between the two delivered signed endpoints? Every candidate must be fully expanded and every group bounded before any claimed saving.
\item Can the positive one-square coprimality mechanism lower degree without losing its operation savings? The present transport majorants are deliberately explicit and prove positivity even when the decoded transport quotient is zero.
\item Can the nonzero-center two-root mechanism improve the signed carrier while keeping its auxiliary root-integrality and fixed-strong-pair interpolation order? In particular a new modulus must not be chosen after the pair it is supposed to constrain.
\item Can the ordinary calibration and each new arithmetic interface be assembled into one end-to-end kernel-checked endpoint? That is distinct unfinished formalization work, not a property implied by the supplied finite checks.
\end{enumerate}
The useful stopping criterion for another optimization round is a fully specified polynomial, an unbounded equivalence proof with explicit dependencies, an authentic coefficient tuple, exact metric evidence and a portable replay. Fewer symbols in an overview, a loose upper bound, or a favorable numerical sample is not such a result.

\end{document}
